% ============================================================
%
%   WORKING NOTES ON A-INFINITY CHIRAL ALGEBRAS
%   AND 3D HOLOMORPHIC-TOPOLOGICAL QFT
%
%   Raeez Lorgat, 2024--2026
%
%   Compile standalone: pdflatex working_notes.tex
%
% ============================================================

\documentclass[11pt]{article}

\usepackage[T1]{fontenc}
\usepackage[utf8]{inputenc}
\usepackage{lmodern}
\frenchspacing

\usepackage[cmintegrals,cmbraces,noamssymbols]{newtxmath}
\usepackage{ebgaramond}

\usepackage[
  activate={true,nocompatibility}, final,
  tracking=true, kerning=true, spacing=true,
  factor=1100, stretch=10, shrink=10
]{microtype}

\usepackage{mleftright}\mleftright

\let\Bbbk\undefined \let\openbox\undefined
\usepackage{amsmath,amssymb,amsthm}
\usepackage{mathtools}
\usepackage{enumitem}
\usepackage[margin=1.15in]{geometry}
\usepackage[unicode]{hyperref}
\usepackage{tikz-cd}

\setlength{\emergencystretch}{3em}
\raggedbottom

% --- Environments ---
\usepackage{thmtools}
\declaretheoremstyle[
  spaceabove=\topsep, spacebelow=\topsep,
  headfont=\normalfont\scshape,
  notefont=\normalfont\itshape,
  bodyfont=\normalfont,
  postheadspace=0.5em, headpunct={.}
]{notesthm}
\declaretheoremstyle[
  spaceabove=\topsep, spacebelow=\topsep,
  headfont=\normalfont\itshape,
  notefont=\normalfont\itshape,
  bodyfont=\normalfont,
  postheadspace=0.5em, headpunct={.}
]{notesdef}

\declaretheorem[style=notesthm, name=Theorem, numberwithin=section]{theorem}
\declaretheorem[style=notesthm, name=Lemma, sibling=theorem]{lemma}
\declaretheorem[style=notesthm, name=Proposition, sibling=theorem]{proposition}
\declaretheorem[style=notesthm, name=Corollary, sibling=theorem]{corollary}
\declaretheorem[style=notesdef, name=Definition, sibling=theorem]{definition}
\declaretheorem[style=notesdef, name=Remark, sibling=theorem]{remark}
\declaretheorem[style=notesdef, name=Observation, sibling=theorem]{observation}
\declaretheorem[style=notesdef, name=Warning, sibling=theorem]{warning}
\declaretheorem[style=notesdef, name=Example, sibling=theorem]{example}
\declaretheorem[style=notesdef, name=Conjecture, sibling=theorem]{conjecture}
\declaretheorem[style=notesdef, name=Question, sibling=theorem]{question}

% --- Macros ---
\newcommand{\cA}{\mathcal{A}}
\newcommand{\cB}{\mathcal{B}}
\newcommand{\cC}{\mathcal{C}}
\newcommand{\cM}{\mathcal{M}}
\newcommand{\cN}{\mathcal{N}}
\newcommand{\cO}{\mathcal{O}}
\newcommand{\cW}{\mathcal{W}}
\newcommand{\cH}{\mathcal{H}}
\newcommand{\cZ}{\mathcal{Z}}
\newcommand{\cD}{\mathcal{D}}
\newcommand{\cI}{\mathcal{I}}
\newcommand{\cL}{\mathcal{L}}
\newcommand{\fg}{\mathfrak{g}}
\newcommand{\fh}{\mathfrak{h}}
\newcommand{\fn}{\mathfrak{n}}
\newcommand{\fsl}{\mathfrak{sl}}
\newcommand{\barB}{\overline{B}}
\newcommand{\B}{\bar{B}}
\newcommand{\C}{\mathbb{C}}
\newcommand{\Z}{\mathbb{Z}}
\newcommand{\R}{\mathbb{R}}
\newcommand{\Q}{\mathbb{Q}}
\newcommand{\bP}{\mathbb{P}}
\newcommand{\bD}{\mathbb{D}}

\newcommand{\Ass}{\mathsf{Ass}}
\newcommand{\Com}{\mathsf{Com}}
\newcommand{\Lie}{\mathsf{Lie}}
\newcommand{\Eone}{\mathsf{E}_1}
\newcommand{\Etwo}{\mathsf{E}_2}
\newcommand{\Einf}{\mathsf{E}_{\infty}}
\newcommand{\Linf}{\mathsf{L}_{\infty}}
\newcommand{\Ainf}{\mathsf{A}_{\infty}}
\newcommand{\SCchtop}{\mathsf{SC}^{\mathrm{ch,top}}}
\newcommand{\Ran}{\operatorname{Ran}}
\newcommand{\Res}{\operatorname{Res}}
\newcommand{\FM}{\mathrm{FM}}
\newcommand{\MC}{\mathrm{MC}}
\newcommand{\ad}{\operatorname{ad}}
\newcommand{\id}{\mathrm{id}}
\newcommand{\Hom}{\operatorname{Hom}}
\newcommand{\RHom}{\operatorname{RHom}}
\newcommand{\End}{\operatorname{End}}
\newcommand{\Ext}{\operatorname{Ext}}
\newcommand{\HH}{\operatorname{HH}}
\newcommand{\Tr}{\operatorname{Tr}}
\newcommand{\Perf}{\operatorname{Perf}}
\newcommand{\ChirHoch}{\mathrm{ChirHoch}}
\newcommand{\gAmod}{\mathfrak{g}^{\mathrm{mod}}_\cA}
\newcommand{\dzero}{d_0}
\newcommand{\dfib}{d_{\mathrm{fib}}}
\newcommand{\Dtot}{D_{\mathrm{tot}}}
\newcommand{\normord}[1]{:\!#1\!:}
\newcommand{\Zder}{\cZ^{\mathrm{der}}_{\mathrm{ch}}}

\DeclareMathOperator{\gr}{gr}
\DeclareMathOperator{\Spec}{Spec}
\DeclareMathOperator{\Conf}{Conf}

% Separator
\newcommand{\sep}{\bigskip\begin{center}$*$\end{center}\medskip}

\title{\textit{Working Notes on $\Ainf$ Chiral Algebras\\
and 3D Holomorphic-Topological QFT}}
\author{Raeez Lorgat}
\date{2024--2026}

\begin{document}
\maketitle

\begin{quote}
\itshape
Volume I took a $1$-form and followed where it led: into the
bar complex, the five theorems, the shadow tower.
These notes are about what happens when the $1$-form acquires
a second direction, what the bar complex becomes when read in
three dimensions, and why the coproduct contains more physics
than the differential.
\end{quote}

\tableofcontents


% =================================================================
\section{Two directions}
\label{sec:two-directions}
% =================================================================

The bar complex of a chiral algebra $\cA$ on a curve $X$ carries
two structures.

\medskip

The first is the \textbf{differential} $d_{\barB}$.  It extracts
OPE residues at collision divisors in $\FM_k(X)$.  Points
approach each other along the curve.  The propagator
$\eta_{ij} = d\log(z_i - z_j)$ detects each collision and records
its residue.  Arnold's relation kills triple residues, so
$d_{\barB}^2 = 0$ at genus $0$.  This is the holomorphic
direction.  It is everything Volume~I studied.

\medskip

The second is the \textbf{coproduct} $\Delta$.  On the tensor
coalgebra $T^c(s^{-1}\bar\cA)$, deconcatenation is
\[
\Delta[a_1 \,|\, \cdots \,|\, a_n]
\;=\;
\sum_{i=0}^{n}
[a_1 \,|\, \cdots \,|\, a_i]
\;\otimes\;
[a_{i+1} \,|\, \cdots \,|\, a_n].
\]
This is the ordered splitting of a word into two subwords.
It does not depend on the curve, the OPE, or the central
charge.  It depends on the ordering.  Points are separated
along a line, and the line is cut at one of the $n+1$
positions between them.  This is the topological direction.

\medskip

The differential lives on $\FM_k(\C)$.  The coproduct lives on
$\Conf_k(\R)$.  Together: a bar element of degree $k$ is
parametrised by
\[
\FM_k(\C) \;\times\; \Conf_k(\R),
\]
the product of holomorphic and topological configuration spaces.

\sep

This product is the operadic fingerprint of a $3$d
holomorphic-topological QFT on $\C_z \times \R_t$.  Observables
factorise holomorphically in $z$ and associatively in $t$.  The
two-coloured Swiss-cheese operad $\SCchtop$ has operation spaces
$\FM_k(\C) \times \Eone(m)$: the closed colour is the
Fulton--MacPherson compactification (holomorphic collisions); the
open colour is the little intervals operad (topological
splittings).  The bar differential is the closed colour.  The bar
coproduct is the open colour.

The same $1$-form $\eta_{ij}$ generates both structures.
Its collision residues produce the differential.  Its ordering
(which point comes first along the time axis) produces the
coproduct.  One form, two directions, one product.

\sep

At genus $0$, $d_{\barB}^2 = 0$ and $\Delta$ is coassociative
and compatible with $d_{\barB}$: these are the Swiss-cheese
relations without obstruction.  At genus $g \geq 1$, the
curvature $\dfib^{\,2} = \kappa(\cA) \cdot \omega_g$ from
Volume~I infects the product:
\[
m_1^2(a) \;=\; [m_0, a], \qquad m_0 = \kappa(\cA) \cdot \omega_g.
\]
The Swiss-cheese algebra becomes \textit{curved}.  The failure of
formality at genus $g \geq 1$ is the bar construction's single
most important feature.  It is not a defect.  It is $\kappa$.


% =================================================================
\section{The recognition theorem}
\label{sec:recognition}
% =================================================================

A factorisation algebra on $\C \times \R$ satisfying local
constancy along~$\R$ is the same thing as an $\SCchtop$-algebra.

The proof uses the product Weiss topology.  A cover of
$U \times I \subset \C \times \R$ is a collection of products
$\{U_\alpha \times I_\beta\}$ refining the product open cover.
The cosheaf of observables on $\C \times \R$ restricts to a
constructible dg-cosheaf on $\FM_k(\C) \times \Conf_k(\R)$.
Weiss descent identifies the structure maps with the operations
of $\SCchtop$.

\begin{theorem}[Recognition; proved]
\label{thm:recognition-wn}
There is an equivalence of $\infty$-categories between:
\begin{enumerate}[label=\textup{(\roman*)},nosep]
\item logarithmic $\SCchtop$-algebras, and
\item factorisation algebras on $\C \times \R$ that are
  holomorphic in $\C$ and locally constant in~$\R$.
\end{enumerate}
\end{theorem}

\noindent
The product topology is not decorative.  It is the entire content
of the theorem: the two colours of the Swiss-cheese operad arise
because the two factors of $\C \times \R$ carry different
topologies (holomorphic versus locally constant).

\medskip

The recognition theorem converts geometric data (a factorisation
algebra on a $3$-manifold) into algebraic data (an operad
algebra).  Without it, the Swiss-cheese structure would be a
presentation, not a property.  With it, the two-coloured operad
is the geometry.


% =================================================================
\section{One algebra, three objects}
\label{sec:three-objects}
% =================================================================

From a single chiral algebra $\cA$, the bar construction
produces three objects.  They are not the same.

\medskip

\noindent
\textbf{(i)} The \textbf{boundary algebra} $\cA$ itself.  This
lives on the boundary $\{0\} \times \R$ of the half-space
$\C_{\geq 0} \times \R$.  It is a chiral algebra: a vertex
algebra, an $\Einf$-chiral structure.  Its data are the OPE
coefficients $c_n^{ab}$ and the conformal weights.

\medskip

\noindent
\textbf{(ii)} The \textbf{bar coalgebra} $\barB(\cA)$.  This is a
factorisation coalgebra on $\Ran(X)$, carrying both the
differential and the coproduct.  It classifies \textit{twisting
morphisms}: universal couplings between $\cA$ and its Koszul
dual~$\cA^!$.  On the Koszul locus,
$\MC(\Hom(\barB(\cA), \cA^!)) \simeq \Perf(\cA^!)$.  This is
the bar complex's native role: a representability theorem for
deformations.

\medskip

\noindent
\textbf{(iii)} The \textbf{chiral derived center}
\[
\Zder(\cA)
\;:=\;
H^*\bigl(C^\bullet_{\mathrm{ch}}(\cA, \cA),\, \delta\bigr).
\]
This is the universal bulk algebra: it classifies observables
of the $3$d theory that act on the boundary.  At the level of
Hochschild cohomology:
$\Zder(\cA) \simeq \HH^\bullet_{\mathrm{ch}}(\cA)$, the chiral
Hochschild cochains.

\sep

The identification ``bar = bulk'' is the most dangerous false
equation in three-dimensional holomorphic-topological field theory.

The bar coalgebra classifies \textit{couplings} between a boundary
and its dual.  The derived center classifies \textit{operators} that
the boundary does not know it has.  These are different functors
applied to different inputs with different outputs.  In a
sentence: bar-cobar is about twisting morphisms; the derived
center is about the acting-on relation.

\begin{center}
\small
\begin{tabular}{@{}lll@{}}
\textbf{Object} & \textbf{Classifies} & \textbf{Functor} \\
\hline
$\barB(\cA)$ & twisting morphisms $\cA \leadsto \cA^!$
  & bar (coalgebra) \\[3pt]
$\Zder(\cA)$ & bulk operators acting on $\cA$
  & chiral Hochschild (algebra) \\[3pt]
$\Omega(\barB(\cA))$ & nothing new: $\cA$ itself
  & cobar (recovers input)
\end{tabular}
\end{center}

\medskip

The cobar of the bar is $\cA$ (Theorem~B of Volume~I).  The
Verdier dual of the bar is $\barB(\cA^!)$ (Theorem~A).  These
are three different operations, producing three different
objects.


% =================================================================
\section{The Swiss-cheese pentagon}
\label{sec:pentagon}
% =================================================================

The $\SCchtop$ operad has two colours: \textit{closed}
(holomorphic, on $\FM_k(\C)$) and \textit{open} (topological,
on $\Eone(m)$).  The structure maps are:

\begin{center}
\small
\renewcommand{\arraystretch}{1.15}
\begin{tabular}{@{}llll@{}}
\textbf{Map} & \textbf{Inputs} & \textbf{Output}
  & \textbf{Interpretation} \\
\hline
$\mu_{k,0}$ & $k$ closed & closed
  & Gerstenhaber bracket \\[3pt]
$\mu_{0,m}$ & $m$ open & open
  & $\Ainf$ products \\[3pt]
$\mu_{k,m}$ & $k$ closed, $m$ open & open
  & brace operations \\[3pt]
$\mu_{k,0}^{\mathrm{oc}}$ & $k$ open & closed
  & \textbf{EMPTY} \\[3pt]
$\mu_{k,m}^{\mathrm{mix}}$ & mixed & mixed
  & SC composition
\end{tabular}
\end{center}

\medskip

The fourth row is the content.  No open inputs produce closed
outputs.  This is not a convention but a structural constraint:
the open colour lives on the boundary $\R$, the closed colour in
the bulk $\C \times \R$, and bulk interactions restrict to
boundaries but not conversely.  Bulk-to-boundary is a restriction
functor; boundary-to-bulk would require extension, which the
holomorphic direction forbids.

The bar complex makes this visible.  The bar \textit{differential}
(closed colour) is the OPE residue extraction: a collision in the
$\C$-plane.  The bar \textit{coproduct} (open colour) is the
deconcatenation: a separation along~$\R$.  Applying the coproduct
and then the differential is a brace operation: splitting a bar
element, then colliding within each piece.  Applying the
differential and then the coproduct is an SC composition:
colliding first, then splitting.  The Stasheff pentagon for the
$\Ainf$ structure and the Gerstenhaber bracket for the closed
colour are the same identity, read on the two factors of
$\FM_k(\C) \times \Eone(m)$.

The Swiss-cheese operad is Koszul.

\begin{theorem}[Homotopy-Koszulity; proved]
\label{thm:homotopy-koszul-wn}
The chiral-topological Swiss-cheese operad $\SCchtop$ is
homotopy-Koszul.
\end{theorem}

\noindent
The proof:
(1)~the classical Swiss-cheese operad is Koszul
(Livernet, Voronov, Ginzburg--Kapranov);
(2)~Kontsevich formality gives a quasi-isomorphism
$\Phi\colon \SCchtop \to \mathrm{SC}$;
(3)~bar-cobar preserves quasi-isomorphisms; two-out-of-three.

The consequences are unconditional.  Bar-cobar on
$\SCchtop$-algebras is a Quillen equivalence, not merely an
adjunction.  Lines are $\cA^!$-modules.  The dg-shifted Yangian
is defined without hypothesis.


% =================================================================
\section{Descent to the classical shadow}
\label{sec:pva-descent}
% =================================================================

The holomorphic Swiss-cheese algebra, projected to cohomology,
produces a Poisson vertex algebra (PVA).  The descent is
controlled by six axioms D1--D6, each proved by a geometric
argument on $\FM_k(\C)$.

\begin{enumerate}[label=D\arabic*.,nosep]
\item \textit{Leibniz rule}:
  $\{a_\lambda bc\} = \{a_\lambda b\}c + b\{a_\lambda c\}$.
\item \textit{Sesquilinearity}:
  $\{\partial a_\lambda b\} = -\lambda\{a_\lambda b\}$.
\item \textit{Skew-symmetry}:
  $\{b_\lambda a\} = -\{a_{-\lambda-\partial} b\}$.
\item \textit{Jacobi identity}:
  $\{a_\lambda\{b_\mu c\}\} - \{b_\mu\{a_\lambda c\}\}
  = \{\{a_\lambda b\}_{\lambda+\mu} c\}$.
\item \textit{Commutativity}:
  $ab = ba$ on cohomology.
\item \textit{Unit}:
  $\{a_\lambda 1\} = 0$.
\end{enumerate}

D2--D5 are proved by the exchange cylinder and the three-face
Stokes identity on $\FM_3(\C)$.  The exchange cylinder is the
bordism
$\FM_2(\C) \times [0,1] \to \FM_3(\C)$
that swaps two points while a third observes; the three-face
identity
\[
d\omega_3 \;=\; \omega_{12,3} - \omega_{13,2} + \omega_{1,23}
\]
on $\partial\FM_3(\C)$ gives skew-symmetry and Jacobi
simultaneously.  D6 uses the factorisation of $H_3(a)$ and the
contractibility of the disc.

The bridge operates in both directions.  F4 (operad $\Rightarrow$
axioms): the $\SCchtop$ structure maps project to PVA operations.
F5 (axioms $\Rightarrow$ operad): a PVA satisfying D1--D6 lifts
to an $\SCchtop$-algebra.  The lift exists because $\SCchtop$ is
homotopy-Koszul: the cobar of the Koszul dual cooperad provides
a cofibrant resolution, and the PVA data define a map from this
resolution.

\sep

The PVA on cohomology $H^\bullet(\cA, Q)$ is $(-1)$-shifted:
the $\lambda$-bracket has shifted parity relative to classical
PVA conventions.  This is not an anomaly.  It is the degree shift
in the desuspension $s^{-1}$ of the bar construction.


% =================================================================
\section{The coproduct}
\label{sec:coproduct}
% =================================================================

In gauge theory (Chern--Simons with boundary algebra
$\widehat\fg_k$), the gluon $J^a$ has a primitive coproduct:
\[
\Delta_z(J^a) \;=\; J^a(w{+}z) \otimes 1
\;+\; 1 \otimes J^a(w).
\]
But the Sugawara composite
$T_{\mathrm{aff}} = \frac{1}{2(k+h^\vee)}\sum_a
\normord{J^a J^a}$ has a nontrivial coproduct.  The cross term
\[
\Delta_z(T_{\mathrm{aff}})
\;=\;
\tau_z(T_{\mathrm{aff}}) \otimes 1
\;+\; 1 \otimes T_{\mathrm{aff}}
\;+\;
\frac{1}{k+h^\vee}\,\Omega(w{+}z,\, w),
\]
where $\Omega = \sum_a J^a \otimes J^a$ is the split Casimir,
is a production vertex: a gluon bound state can dissociate into
its constituents.  The gluon IS the connection.  It splits
because the connection is a primitive field.

\sep

In gravity (Virasoro boundary), the stress tensor $T$ is the
graviton.  The coproduct is primitive at every arity:
\begin{equation}\label{eq:gravity-primitive-wn}
\Delta_z^W(x) \;=\; \tau_z(x) \otimes 1 \;+\; 1 \otimes x
\end{equation}
for every element $x$ of the $\cW$-algebra, with no corrections
whatsoever.

The mechanism is a ghost-number obstruction intrinsic to
Drinfeld--Sokolov reduction; the proof is given in
\S\ref{sec:primitive-coproduct}
(Theorem~\ref{thm:gravitational-primitivity-wn}).

In gauge theory, the gluon IS the connection.  In gravity, the
graviton IS the curvature.  A connection splits (it is the
primitive field of the theory).  A curvature does not (it is a
composite, and the ghost-number filtration forbids the splitting).
This is the algebraic content of the statement that gravity does
not produce particles; it only braids them.

\begin{center}
\small
\renewcommand{\arraystretch}{1.15}
\begin{tabular}{@{}llll@{}}
 & \textbf{Chern--Simons} & \textbf{Gravity} \\
\hline
Boundary algebra & $\widehat\fg_k$ & $\mathrm{Vir}_c$ \\[3pt]
Primitive field & $J^a$ (connection) & $T$ (curvature) \\[3pt]
$\Ainf$ tower & formal ($m_k^H = 0$ for $k \ge 3$, class L)
  & infinite ($m_k \neq 0\ \forall k$, class M) \\[3pt]
Coproduct $\Delta_z$ & nontrivial (Casimir cross-term)
  & primitive (ghost obstruction) \\[3pt]
$r$-matrix & $\Omega_\fg/z$ & $(c/2)/z^3 + 2T/z$ \\[3pt]
Shadow depth & $3$ & $\infty$
\end{tabular}
\end{center}

\medskip

All nonlinearity of gravity resides in the $\Ainf$ products
$\{m_n^W\}_{n \geq 3}$ and in the collision residue
$r^{\mathrm{Vir}}(z) = (c/2)/z^3 + 2T/z$.  The coproduct carries
no information.  This is the central algebraic fact about
three-dimensional quantum gravity.


% =================================================================
\section{Five worked examples}
\label{sec:five-examples}
% =================================================================

\subsection*{Chern--Simons}

Gauge group $G$, level $k$.  Boundary algebra:
$\widehat\fg_k$.  Class L, shadow depth $3$.
$r$-matrix: $r(z) = \Omega_\fg/z$ (Casimir, single pole).
Coproduct: nontrivial (Casimir cross-term on composites).
$\kappa = \dim(\fg)(k+h^\vee)/(2h^\vee)$.
Complementarity: $\kappa + \kappa^! = 0$ by Feigin--Frenkel.

The entire $\Ainf$ structure is formal: $m_k = 0$ for $k \geq 3$
on cohomology.  The bar complex is quadratic.  The Koszul dual
$\widehat\fg_{-k-2h^\vee}$ is another affine algebra at the
reflected level.  Lines: highest-weight modules for the dual.
On the integrable locus ($k \in \Z_{\geq 0}$), the line category
is $\mathrm{Rep}_q(\fg)$ with $q = e^{\pi i/(k+h^\vee)}$,
recoverable by the Kazhdan--Lusztig equivalence.

\subsection*{Gravity}

Boundary algebra: $\mathrm{Vir}_c$.  Class M, shadow depth
$\infty$.  $r$-matrix: $r(z) = (c/2)/z^3 + 2T/z$ (two poles,
quasi-trigonometric).  Coproduct: primitive.
$\kappa = c/2$; $\kappa^! = (26-c)/2$;
$\kappa + \kappa^! = 13 \neq 0$.

The quartic OPE pole forces an infinite $\Ainf$ tower.  Every
$m_k$ is nonzero for generic $c$, computed by a wheel diagram on
$K_k$ with $k-1$ cubic vertices.  The quartic contact invariant
$Q^{\mathrm{ct}}_{\mathrm{Vir}} = 10/[c(5c+22)]$
has poles at $c = 0$ and $c = -22/5$: the Kac determinant zeros
at weight~$4$.  Self-duality at $c = 13$, not $c = 26$.

\subsection*{Twisted holography}

Costello--Gaiotto.  Boundary algebra: $V_k(\fg)$ (affine
Kac--Moody).  Class L, shadow depth $3$.  Coproduct: nontrivial.
The holographic dictionary is the Koszul triangle read through
Volume~I's five theorems: bar complex exists (Theorem~A), boundary
reconstructs bulk (Theorem~B), Neumann/Dirichlet decomposition
(Theorem~C), genus tower by $\kappa$ (Theorem~D), BV origin of
Hochschild complex (Theorem~H).

\subsection*{M2-brane}

From the $\cN = 6$ ABJM theory.  Boundary algebra: supersymmetric
current algebra.  Coproduct: nontrivial (fundamental fields split).
The M2 example has no DS reduction: the $\cW$-algebra structure is
absent, and the coproduct obstruction does not apply.  This is a
gauge theory, not a gravity theory, and the coproduct confirms it.

\subsection*{M5-brane / 6d $(2,0)$}

From the $\cN = (2,0)$ theory on $\C \times \R \times S^3$.
Boundary algebra: $\cW_N$ at large $N$.  Class M, shadow depth
$\infty$.  Coproduct: primitive (by the same ghost-number
argument as Virasoro, since $\cW_N$ arises by DS reduction).

\begin{center}
\small
\renewcommand{\arraystretch}{1.15}
\begin{tabular}{@{}lccccc@{}}
 & CS & Gravity & Twisted hol. & M2 & M5 \\
\hline
Class & L & M & L & L & M \\
$r_{\max}$ & $3$ & $\infty$ & $3$ & $3$ & $\infty$ \\
$\Delta_z$ & nontrivial & primitive & nontrivial
  & nontrivial & primitive \\
DS origin & no & yes & no & no & yes \\
$\kappa + \kappa^!$ & $0$ & $13$ & $0$ & $0$ & $\varrho K_N$
\end{tabular}
\end{center}

\medskip

The dichotomy is clean: DS reduction produces primitive
coproducts and infinite shadow depth.  Its absence permits
nontrivial coproducts and finite termination.


% =================================================================
\section{The annulus}
\label{sec:annulus}
% =================================================================

Cut the circle $S^1$ by two arcs.  Each arc is an interval
$I \cong \R$.  The factorisation homology of the open-sector
category $\cC_{\mathrm{op}}$ on $S^1$ is computed by excision:
\[
\int_{S^1} \cC_{\mathrm{op}}
\;\simeq\;
\cC_{\mathrm{op}}(I_1)
\;\otimes_{\cC_{\mathrm{op}}(I_1 \cap I_2)}^{\mathbb{L}}\;
\cC_{\mathrm{op}}(I_2).
\]
The two-arc cover produces the two-sided bar complex, which is
the cyclic bar complex.  In the homotopy category:
\begin{equation}\label{eq:annulus-trace-wn}
\int_{S^1} \cC_{\mathrm{op}}
\;\simeq\;
\HH_*(\cC_{\mathrm{op}}).
\end{equation}

This is the annulus trace.  The annulus $S^1 \times [0,1]$ is the
simplest bordered surface with one boundary component.  Its
factorisation homology is the Hochschild homology of the
open-sector category.  The trace is the clutching operation: it
restores the tangential basepoint $\tau_p$ to close the interval
$I_p$ back to a circle.

\sep

For affine Kac--Moody at integrable level $k$:
$\HH_0(\cC_{\mathrm{op}})$ is spanned by the characters
$\chi_j(\tau) = \mathrm{tr}_{L_j}(q^{L_0 - c/24})$
for $j = 0, 1, \ldots, k$.  The modular $S$-matrix acts on
this space.  Modular invariance of the partition function
$Z(\tau) = \sum_j |\chi_j|^2$ is the statement that the trace
is $SL(2,\Z)$-equivariant.

\medskip

The annulus trace is the \textbf{first modular shadow of the
open sector}.  It is the genus-$1$ datum, the simplest case of
the trace-and-clutching construction that produces genus-$g$
amplitudes from the open category.  Modularity is not an axiom
on the closed algebra.  It is a property of the trace on the
open sector.


% =================================================================
\section{The line category}
\label{sec:line-category}
% =================================================================

On the chirally Koszul locus:
\begin{equation}\label{eq:lines-koszul-wn}
\cC_{\mathrm{line}} \;\simeq\; \cA^!\text{-}\mathbf{mod}.
\end{equation}
Line operators are modules for the Koszul dual.

The proof is Koszul duality on the open colour.  The bar complex
$\barB(\cA)$ is a dg coalgebra; its linear dual
$\barB(\cA)^\vee = \barB(\cA^!)$ (Verdier duality, Theorem~A
of Volume~I) is a dg algebra.  Maurer--Cartan deformations of
the universal twisting morphism
$\alpha\colon \barB(\cA) \to \cA^!$ are in bijection with
$\cA^!$-module structures:
\[
\MC\bigl(\Hom(\barB(\cA),\, \cA^!)\bigr)
\;\simeq\;
\Perf(\cA^!).
\]
A line is not a module that happens to be placed along $\R$.
A module is a line: the representation category of $\cA^!$ IS the
category of boundary conditions for the theory restricted
to~$\R$.

\sep

The meromorphic tensor structure on $\cC_{\mathrm{line}}$ comes
from the $r$-matrix.  Two line operators at spectral
parameters $z_1$ and $z_2$ tensor along the open direction with
braiding
\[
R(z_1 - z_2)\colon M_1 \otimes M_2
\;\xrightarrow{\sim}\;
M_2 \otimes M_1,
\]
where $R(z)$ is the quantum $R$-matrix built from the bulk-boundary
composition (not from a quantum group, though they agree on the
evaluation locus).  For Chern--Simons:
$R(z) = 1 + \hbar\,\Omega_\fg/z + O(\hbar^2)$.  For gravity:
$R(z) = 1 + \hbar\,((c/2)/z^3 + 2T/z) + O(\hbar^2)$.

The $r$-matrix is the genus-$0$, arity-$2$, ordered shadow of
$\Theta_\cA$:
\[
r(z) \;=\; \Res^{\mathrm{coll}}_{0,2}(\Theta_\cA).
\]
Volume~I computes the collision residue.  Volume~II reads the
result as a spectral braiding on lines.


% =================================================================
\section{The benchmark}
\label{sec:benchmark}
% =================================================================

$V_k(\fsl_2)$ on $(\bP^1, \{0, \infty\})$.  This is the
simplest non-trivial example: one curve with two punctures, one
simple Lie algebra of rank $1$, one level parameter $k$.  Every
object in the platonic datum can be computed.

\medskip

\noindent
\textbf{Tangential log curve.}
$(X, D, \tau) = (\bP^1, \{0,\infty\}, \{\tau_0, \tau_\infty\})$
with tangential basepoints in the positive-real direction.  The
boundary intervals $I_0, I_\infty \cong \R$ are the time axes at
each puncture.

\medskip

\noindent
\textbf{Open sector.}
$\cC_{\mathrm{op}} = V_k(\fsl_2)\text{-mod}$ at each puncture.
At integrable level $k$, this is $\mathrm{Rep}_q(\fsl_2)$ with
$q = e^{i\pi/(k+2)}$, a semisimple category with $k+1$ simple
objects $L_0, L_1, \ldots, L_k$.

\medskip

\noindent
\textbf{Boundary algebra.}
$A_b = V_k(\fsl_2)$.  The $\lambda$-bracket:
$\{J^a{}_\lambda J^b\} = f^{ab}{}_c J^c + k\kappa^{ab}\lambda$.

\medskip

\noindent
\textbf{Koszul dual.}
$V_k(\fsl_2)^! = V_{-k-4}(\fsl_2)$.  The Feigin--Frenkel
involution $k \mapsto -k - 2h^\vee = -k - 4$.

\medskip

\noindent
\textbf{$r$-matrix.}
$r(z) = \Omega_{\fsl_2}/z = (J^a \otimes J_a)/z$, a single-pole
solution of the classical Yang--Baxter equation.

\medskip

\noindent
\textbf{Modular characteristic.}
$\kappa = 3(k+2)/4$.  Complementarity: $\kappa + \kappa^! = 0$.

\medskip

\noindent
\textbf{Derived center.}
$\Zder(V_k(\fsl_2)) \simeq \HH^\bullet_{\mathrm{ch}}(V_k(\fsl_2))$.
At the critical level $k = -2$: functions on $\fsl_2$-opers.

\medskip

\noindent
\textbf{Line category.}
$\cC_{\mathrm{line}} \simeq V_{-k-4}(\fsl_2)\text{-mod}$.
At integrable levels: $\mathrm{Rep}_{q^{-1}}(\fsl_2)$.

\medskip

\noindent
\textbf{Annulus trace.}
$\HH_0(\cC_{\mathrm{op}}) = \bigoplus_{j=0}^k \C \cdot \chi_j$.
The modular $S$-matrix:
\[
S_{jj'} \;=\; \sqrt{\frac{2}{k+2}}\,
\sin\!\left(\frac{\pi(j+1)(j'+1)}{k+2}\right).
\]
At $k = 1$: $S = \frac{1}{\sqrt{2}}
\begin{psmallmatrix} 1 & 1 \\ 1 & -1 \end{psmallmatrix}$,
the Hadamard matrix.

\medskip

\noindent
\textbf{Factorisation homology.}
$\int_{\C^\times}\!\cC_{\mathrm{op}}
\simeq \mathrm{CB}_{0,2}(V_k(\fsl_2))$:
genus-$0$ conformal blocks with two insertions.  For simple
modules $L_j$ and $L_{j'}$:
$\dim \mathrm{CB}_{0,2}(L_j, L_{j'}) = \delta_{j,j'}$.

\medskip

Every entry is proved.  The benchmark is the frame atom of
Volume~II: simple enough to compute completely, rich enough to
exhibit every feature of the theory.


% =================================================================
\section{Modularity last}
\label{sec:modularity-last}
% =================================================================

The standard approach to modular invariance begins with the
closed-string partition function and imposes modular invariance
as an axiom.  The Swiss-cheese framework reverses this.

Modularity is a derived property of the trace on the open sector
(\S\ref{sec:annulus}).

\sep

The modular Swiss-cheese operad extends $\SCchtop$ by replacing
trees with stable graphs carrying two edge colours: closed
edges (genus direction, carrying $\hbar$-weight, recording
non-separating degenerations) and open edges (boundary direction,
carrying insertion data, recording separations along $\R$).

The relative Feynman transform $F_{\mathrm{rel}}\cO$ sums over
stable graphs with both edge colours.  The differential contracts
one edge; $d^2 = 0$ is the codimension-$2$ boundary cancellation
on the moduli space of bordered Riemann surfaces.

At genus $g \geq 1$, the curvature
$\dfib^{\,2} = \kappa \cdot \omega_g$ forces a passage from
ordinary to curved Swiss-cheese algebras.  The modular MC element
\begin{equation}\label{eq:oc-mc-wn}
\Theta^{\mathrm{oc}}_\cA
\;:=\;
\Theta_\cA + \sum_{j} \mu^{\cM_j}
\end{equation}
packages both the bar-intrinsic MC element $\Theta_\cA$ (from
Volume~I) and the $\Ainf$-module structures $\mu^{\cM_j}$ at
punctures.  Its MC equation decomposes into four components:
\begin{alignat*}{2}
&\text{pure closed:} &&\quad d_0\,\Theta
  + \tfrac{1}{2}[\Theta, \Theta] = 0 \\[3pt]
&\text{pure open:} &&\quad \partial\mu
  + \mu \circ \mu = 0 \\[3pt]
&\text{mixed:} &&\quad d_0\,\mu
  + [\Theta, \mu] + \cdots = 0 \\[3pt]
&\text{clutching:} &&\quad \text{(genus recursion)}
\end{alignat*}
The clutching component is the genus recursion from Volume~I.
Each genus is determined by all lower genera.  No parameter
is free.

\sep

The annulus trace (the $S^1$ factorisation homology of
$\cC_{\mathrm{op}}$) produces the genus-$1$ modular datum.
The clutching of $g$ annuli produces the genus-$g$ datum.
Modularity emerges from the composition law of traces, not
from an a priori constraint on the closed algebra.

The programme status:

\begin{center}
\small
\renewcommand{\arraystretch}{1.15}
\begin{tabular}{@{}ll@{}}
\textbf{Result} & \textbf{Status} \\
\hline
Annulus trace $\simeq \HH_*(\cC_{\mathrm{op}})$ & proved \\[3pt]
Full modular cooperad & programme (in progress) \\[3pt]
Modular Swiss-cheese $d^2 = 0$ & proved (by $\partial^2 = 0$ on
  $\overline\cM_{g,n}$) \\[3pt]
Curvature $= \kappa \cdot \omega_g$ & proved (Volume~I)
\end{tabular}
\end{center}


% =================================================================
\section{The gravitational coproduct is primitive}
\label{sec:primitive-coproduct}
% =================================================================

\begin{theorem}[Gravitational primitivity; proved]
\label{thm:gravitational-primitivity-wn}
For every principal Drinfeld--Sokolov reduction of every simple
Lie algebra, the transferred dg-shifted Yangian coproduct on the
$\cW$-algebra is strictly primitive:
\[
\Delta_z^W(x)
\;=\;
\tau_z(x) \otimes 1
\;+\; 1 \otimes x
\]
for every $x$, at every arity, with no corrections.
\end{theorem}

\begin{proof}
Three lines.

The homological perturbation lemma expresses the transferred
coproduct as a sum over decorated trees.  Each tree takes
$\cW$-algebra inputs (ghost number $0$), applies the BRST
homotopy $h$ (which raises or lowers ghost number by $1$), and
outputs a tensor product.  Every tree with at least one internal
vertex produces an output at ghost number $\pm 1$.  The projection
$p \otimes p$ onto ghost-number-$0$ cohomology annihilates both
signs.  The only surviving contribution is the trivial tree
(no internal vertices), which gives the primitive term
$\tau_z(x) \otimes 1 + 1 \otimes x$.
\end{proof}

\sep

The proof works for every principal DS reduction because the
ghost-number obstruction is intrinsic to the BRST structure, not
to the specific Lie algebra.  The ghost field $c \in \fn_+^*$
has ghost number $+1$; the antighost $b \in \fn_+$ has ghost
number $-1$; the BRST homotopy $h$ carries ghost number $\pm 1$;
and the projection $p$ annihilates everything outside ghost
number $0$.

The physical content is the no-particle-production principle of
\S\ref{sec:coproduct}: the graviton is a composite (the stress
tensor, built from the Sugawara construction) and not a fundamental
field, so the ghost-number filtration forbids the splitting at
every order.

For Chern--Simons, the fundamental field $J^a$ has ghost number
$0$, so the BRST homotopy produces ghost-number-$0$ outputs.
The Casimir cross-term survives projection.  The gluon splits.

The asymmetry between gauge theory and gravity is an asymmetry
between fundamental and composite fields, detected by the ghost
number.


% =================================================================
\section{The four-stage architecture}
\label{sec:four-stages}
% =================================================================

The open/closed programme has four stages.  Each depends on its
predecessor.

\medskip

\noindent
\textbf{Stage 1: local one-colour.}
The input is a single chiral algebra $\cA$.  Outputs:

\begin{enumerate}[label=\textup{(\alph*)},nosep]
\item The bar coalgebra $\barB(\cA)$ exists as a factorisation
  coalgebra (Volume~I, Theorem~A).
\item On the Koszul locus, MC couplings
  $\simeq \Perf(\cA^!)$ (Koszul duality on the open colour).
\item The bulk is the chiral derived center:
  $\Zder(\cA) \simeq \HH^\bullet_{\mathrm{ch}}(\cA)$
  (boundary-linear exact theorem).
\item Braces: $C^\bullet_{\mathrm{ch}}(\cA, \cA)$ carries a brace
  dg-algebra structure from the $\SCchtop$ action.
\end{enumerate}

No global factorisation or modularity is used.  Everything here
is the universal property of the bar construction and the
homotopy-Koszulity of $\SCchtop$.  Proved.

\medskip

\noindent
\textbf{Stage 2: the open primitive.}
The primitive datum is the open/closed factorisation dg-category
$\cC$ on a pointed curve with tangential basepoints
$(X, D, \tau)$.

\begin{enumerate}[label=\textup{(\alph*)},nosep]
\item For any compact generator $b \in \cC(J)$, the boundary
  algebra $A_b = \RHom_{\cC(J)}(b,b)$ is an $\Ainf$ chiral
  algebra.
\item Different compact generators produce Morita-equivalent
  boundary algebras: the category $\cC$ is the invariant,
  $A_b$ is a chart.
\item The bar complex $\barB(A_b)$ encodes the twisting data
  of $\cC$, and bar-cobar is the representability theorem for
  deformations of $b$ inside $\cC$.
\end{enumerate}

Proved.

\medskip

\noindent
\textbf{Stage 3: globalise.}
Replace $(X, D, \tau) = (\C, \{0\}, \tau_0)$ by a general
tangential log curve.  The real oriented blowup
$\pi\colon \widetilde{X}_D \to X$ replaces each puncture $p_i$
by a circle $S^1_{p_i}$; the tangential basepoint $\tau_{p_i}$
cuts this circle into an interval $I_{p_i} \cong \R$.

The mixed configuration spaces
$\Conf^{\mathrm{oc}}_{I,\mathbf{m}}(X, D, \tau)
= \Conf_I(X \setminus D) \times \prod_p \Conf^<_{m_p}(\R)$
parametrise bulk insertions and boundary insertions together.
The mixed Ran space is their colimit.

Local-global bridge: the local one-colour theorems of Stage 1
extend to global statements via descent along the mixed Ran space.
Proved locally, the global extension requires compact generation
and derived-center hypotheses.

\medskip

\noindent
\textbf{Stage 4: modularise.}
The trace operation on $\cC_{\mathrm{op}}$ (closing an interval
to a circle) produces genus-$g$ data from genus-$0$ data.

\begin{enumerate}[label=\textup{(\alph*)},nosep]
\item The annulus trace
  $\mathrm{Tr}_\cA \simeq \HH_*(\cC_{\mathrm{op}})$ is the
  first modular shadow.  Proved.
\item The full modular cooperad structure on the trace (the
  composition law for genus-$g$ amplitudes via clutching of $g$
  handles) is a programme.
\end{enumerate}

\sep

The architecture is sequential.  Stage $k$ depends on Stage $k-1$.
One should never discuss Stage 4 before establishing Stage 1.
The order is: objects first, then categories, then globalisation,
then modularity.


% =================================================================
\section{The Koszul triangle}
\label{sec:koszul-triangle}
% =================================================================

On the chirally Koszul locus, the corrected bulk-boundary-line
triangle is:
\begin{equation}\label{eq:triangle-wn}
A_{\mathrm{bulk}}
\;\simeq\;
\Zder(B_\partial)
\;\simeq\;
\Zder(\cC_{\mathrm{line}})
\;\simeq\;
\HH^\bullet_{\mathrm{ch}}(\cA^!),
\qquad
\cC_{\mathrm{line}} \simeq \cA^!\text{-}\mathbf{mod}.
\end{equation}

Three vertices:
\begin{itemize}[nosep]
\item The \textbf{bulk} is
  $A_{\mathrm{bulk}} \simeq \cO(T^*[-1]\cL)$: polyvector fields
  on the Lagrangian.  Note: the restriction
  $\cO(\cM)|_\cL = \cO(\cL)$ is the boundary, not the bulk.
  The bulk retains the full $(-1)$-shifted cotangent fibre.
\item The \textbf{boundary} is $B_\partial = \cO(\cL)$: functions
  on the Lagrangian.
\item The \textbf{line category} is
  $\cC_{\mathrm{line}} = \mathfrak{S}_b\text{-mod}$: modules for
  the self-intersection algebra
  $\mathfrak{S}_b = \cL \times_\cM \cL$.
\end{itemize}

The equivalence $\Zder(B_\partial) \simeq A_{\mathrm{bulk}}$ is
the Hochschild--Kostant--Rosenberg theorem: Hochschild cochains of
$\cO(\cL)$ are polyvector fields on $T^*[-1]\cL \simeq \cM|_\cL$,
by the Lagrangian condition
$T_\cL\cM / T\cL \simeq T^*\cL[-1]$.

The equivalence $\cC_{\mathrm{line}} \simeq \cA^!\text{-mod}$ is
the convolution-module model: the self-intersection algebra on
$\mathfrak{S}_b$ is the Koszul dual, and its modules are lines.

Both equivalences require the Lagrangian condition.  This is the
geometric origin of the algebraic fact that the boundary is NOT
the bulk and lines are NOT bulk modules.

\sep

The bar complex presents the Swiss-cheese algebra, as the
Steinberg variety presents the Hecke algebra.

This is not analogy.  Both are algebro-geometric presentations of
operadic structures via configuration-space resolutions.  The
Steinberg variety $\mathfrak{S}_b = \cL \times_\cM \cL$ is the
self-intersection of the Lagrangian $\cL$ in the moduli stack
$\cM$; its convolution algebra is the Hecke algebra.  The bar
complex $\barB(\cA) = T^c(s^{-1}\bar\cA)$ with differential from
collision residues and coproduct from deconcatenation is the
Swiss-cheese algebra on $\FM_k(\C) \times \Conf_k(\R)$; its
convolution with the Koszul dual is the line category.  The
two become the same construction at the level of factorisation
homology.


% =================================================================
\section{What we do not know}
\label{sec:what-we-do-not-know}
% =================================================================

\subsection*{Three falsification tests}

\textbf{Test 1.}
The global bulk-boundary-line triangle
(equation~\eqref{eq:triangle-wn}) holds unconditionally in the
boundary-linear exact sector.  Whether it extends to all HT
theories depends on two unverified hypotheses: compact generation
of the open-sector category, and the derived center being a
quasi-isomorphism to the bulk.  These hypotheses are verified for
Chern--Simons and for inputs from the Costello--Davison--Gaiotto
framework.  They are not verified for gravity.  A counterexample
would mean that the bulk of a gravitational theory is not the
derived center of its boundary: the simplest open/closed
identification would fail.

\medskip

\textbf{Test 2.}
The gravitational primitivity theorem
(\S\ref{sec:primitive-coproduct}) is proved for principal DS
reductions.  For non-principal DS reductions (hook-type in type A,
arbitrary nilpotent $f$ in general), the ghost-number argument
applies verbatim, but the coproduct on the BRST cohomology has
not been computed.  Hook-type in type A is the proved corridor.
Full non-principal primitivity is conjectural.

\medskip

\textbf{Test 3.}
The modular cooperad structure on the annulus trace (Stage 4 of
the architecture) is a programme, not a theorem.  The annulus
trace itself is proved: $\mathrm{Tr}_\cA \simeq
\HH_*(\cC_{\mathrm{op}})$.  The full cooperad structure (clutching
laws for genus $g \geq 2$, compatibility with the modular operad
on $\overline\cM_{g,n}$) requires the bordered FM compactification
and its log-FM extension, which is the frontier of Mok's
programme.

\subsection*{Programme status ledger}

\begin{center}
\small
\renewcommand{\arraystretch}{1.15}
\begin{tabular}{@{}llc@{}}
\textbf{Component} & \textbf{What it requires} & \textbf{Status} \\
\hline
$\SCchtop$ homotopy-Koszulity & Kontsevich + Livernet &
  \textsc{green} \\[3pt]
Recognition theorem & Product Weiss descent &
  \textsc{green} \\[3pt]
PVA descent D1--D6 & FM Stokes calculus &
  \textsc{green} \\[3pt]
Gravitational primitivity & Ghost-number obstruction &
  \textsc{green} \\[3pt]
Bulk $=$ derived center (boundary-linear) &
  Stage 1 local theorem & \textsc{green} \\[3pt]
Bulk $=$ derived center (global) &
  Compact generation & \textsc{amber} \\[3pt]
Lines $= \cA^!$-mod & Koszul locus &
  \textsc{green} \\[3pt]
Annulus trace & $S^1$ excision &
  \textsc{green} \\[3pt]
Full modular cooperad & Bordered FM + log extension &
  \textsc{amber} \\[3pt]
BV/BRST $=$ bar at $g \geq 1$ & Genus-$g$ QME
  & \textsc{red}
\end{tabular}
\end{center}

\subsection*{Five false ideas to dismiss}

\begin{enumerate}[nosep]
\item ``The bar complex is the bulk algebra.''  The bar complex
  classifies twisting morphisms.  The bulk is the derived center.
  These are different objects produced by different functors.
\item ``The boundary algebra determines the theory.''  The boundary
  algebra is a chart of the open-sector category, not its essence.
  Different compact generators produce Morita-equivalent charts.
  The category is the invariant.
\item ``Modularity is a property of the closed algebra.''  It is a
  property of the trace on the open sector.  The closed algebra
  inherits modularity from the clutching operation, not the reverse.
\item ``Gravity and gauge theory differ only in the central
  charge.''  They differ in the coproduct.  Gauge theory has
  nontrivial $\Delta_z$; gravity has primitive $\Delta_z$.  The
  ghost-number obstruction is structural, not parametric.
\item ``The Koszul dual lives in the bulk.''  It lives on the
  boundary: $\cA^!$ is the boundary algebra for line operators,
  not a bulk observable.
\end{enumerate}


% =================================================================
\section{The completed platonic datum}
\label{sec:platonic}
% =================================================================

From a single chiral algebra $\cA$, the completed platonic datum
assembles eight components:
\[
\Pi^{\mathrm{oc}}_X(\cA)
\;=\;
\bigl(\barB_X,\;\; \Omega_X,\;\; \Theta_\cA,\;\;
\Zder,\;\; \mathrm{Tr}_\cA,\;\;
\HH_*,\;\; \mathfrak{R}^{\mathrm{oc}}_\bullet,\;\;
\nabla^{\mathrm{hol}}\bigr).
\]

\begin{center}
\small
\renewcommand{\arraystretch}{1.15}
\begin{tabular}{@{}lll@{}}
\textbf{Component} & \textbf{Object} & \textbf{Role} \\
\hline
$\barB_X$ & bar coalgebra & classifies twisting morphisms \\[3pt]
$\Omega_X$ & cobar algebra & recovers $\cA$ (Theorem B) \\[3pt]
$\Theta_\cA$ & universal MC class & organises nonlinear tower \\[3pt]
$\Zder$ & chiral derived center & universal bulk \\[3pt]
$\mathrm{Tr}_\cA$ & annulus trace
  & first modular shadow \\[3pt]
$\HH_*$ & Hochschild homology & factorisation on $S^1$ \\[3pt]
$\mathfrak{R}^{\mathrm{oc}}_\bullet$ & nonlinear resonance tower
  & quartic and higher obstructions \\[3pt]
$\nabla^{\mathrm{hol}}$ & holographic connection
  & genus-$g$ transport
\end{tabular}
\end{center}

\medskip

The first six components are proved.  The nonlinear resonance
tower is proved at finite arity (Volume~I, shadow Postnikov tower).
The holographic connection is defined but its full modular
extension is the frontier.

\sep

The completed platonic datum is the constructive endpoint of the
theory.  Given any chiral algebra, compute $\Pi^{\mathrm{oc}}_X$.
Theorems A--H of Volume~I are structural properties of
$\Pi^{\mathrm{oc}}_X$.  The five worked examples of
\S\ref{sec:five-examples} are instances of this computation.

The relationship between the six-fold platonic package $\Pi_X(L)$
of Volume~I (built from a Lie conformal algebra $L$) and the
eight-fold $\Pi^{\mathrm{oc}}_X(\cA)$ (built from a chiral
algebra $\cA$) is the factorisation envelope: the functor
$L \mapsto U_X(L)$ that produces the enveloping vertex algebra
from $L$.  The six-fold package lives at the input; the
eight-fold package lives at the output.  The two additional
components (derived center and annulus trace) are intrinsically
$3$-dimensional: they require the second direction.


% =================================================================
\section{From one quartic pole}
\label{sec:one-quartic-pole}
% =================================================================

The Virasoro $\lambda$-bracket
\[
\{T{}_\lambda T\}
\;=\;
\partial T + 2T\lambda + \frac{c}{12}\,\lambda^3
\]
has three terms: the simple pole ($\partial T$, the derivative),
the double pole ($2T\lambda$, the stress tensor), and the
quartic pole ($\frac{c}{12}\lambda^3$, the central charge).

The quartic pole is the engine.  It forces non-associativity of
$m_2$.  Non-associativity of $m_2$ forces $m_3$ by the Stasheff
identity $\partial m_3 + m_2 \circ m_2 = 0$.  Each $m_k$ forces
$m_{k+1}$.  The $\Ainf$ structure is genuinely infinite (class~M,
shadow depth $\infty$).

The collision residue of the $\lambda$-bracket is the gravitational
$r$-matrix:
\[
r(z) \;=\; (c/2)/z^3 + 2T/z.
\]
Two poles: a triple pole from the central charge, a simple pole from
the stress tensor.  The pole orders are one less than the OPE (the
$d\log$ kernel absorbs one power).  The $r$-matrix has no
even-order poles for a bosonic algebra.

\sep

From this single input:
\begin{enumerate}[label=\textup{(\Roman*)},nosep]
\item The $\Ainf$ tower: $m_3, m_4, \ldots$, all nonzero for
  generic $c$.  The quartic contact invariant
  $Q^{\mathrm{ct}} = 10/[c(5c+22)]$.
\item The Koszul dual: $\mathrm{Vir}_{26-c}$.  Self-duality at
  $c = 13$.  The depth-zero resonance shadow.
\item The genus tower:
  $\sum_{g \geq 1} F_g\,\hbar^{2g}
  = (c/2)(\hat{A}(i\hbar) - 1)$,
  with $\kappa_{\mathrm{eff}} = (c-26)/2$.
\item The Koszul triangle: boundary $\mathrm{Vir}_c$, bulk
  $\C[\![c]\!]$, lines modelled by $\mathrm{Vir}_{26-c}$-modules.
\item The gravitational coproduct: primitive (ghost-number
  obstruction from DS reduction of $\widehat\fg_k$).
\end{enumerate}

\medskip

The transgression algebra of Volume~I's working notes completes
the picture.  The Koszul dual $B = \mathrm{Vir}_{26-c}$ has
curvature $\lambda = \kappa(B) \cdot \omega_1
= \frac{26-c}{2}\,\omega_1$.  The transgression algebra
\[
B_\lambda \;=\; B\langle\eta\rangle\,\big/\,(\eta^2 - \lambda),
\qquad |\eta| = 1,
\]
adjoins a Clifford generator.  The secondary anomaly
$u = \eta^2 = \lambda$ is linear in $\kappa$ (a class, not a
scalar).  The gravity dichotomy:

\begin{center}
\small
\renewcommand{\arraystretch}{1.15}
\begin{tabular}{@{}lcc@{}}
& $c \neq 26$ & $c = 26$ \\
\hline
$u$ & $\neq 0$ & $= 0$ \\[3pt]
Clifford algebra & matrix ($2^g$-dim) & exterior ($2^{2g}$-dim)\\[3pt]
Regime & perturbative & topological \\[3pt]
Genus tower & resummable ($\hat{A}$-genus) & collapsed \\[3pt]
Surface topology & Morita-trivial & genuine
\end{tabular}
\end{center}

\medskip

The MC recursion determines the genus-$g$ partition function from
the genus-$0$ $\lambda$-bracket.  No sum over topologies.  No free
parameters.  The Maloney--Witten problem of enumeration (which
topologies to include) is replaced by a problem of algebra (which
Maurer--Cartan equation to solve).

\begin{question}
Does the graded character
$\hat{Z}_1^{\mathrm{grav}}(c;\,\tau)
:= \chi_{\mathrm{gr}}(\mathfrak{G}_1^{\mathrm{grav}}(c);\,q)$
have non-negative Fourier coefficients for all $c > 1$?
\end{question}


% =================================================================
\section{The Steinberg principle}
\label{sec:steinberg}
% =================================================================

The parallel announced in \S\ref{sec:koszul-triangle} is not
analogy; it is a theorem.

The Steinberg variety $\mathfrak{S} = \widetilde\cN
\times_\cN \widetilde\cN$ is the self-intersection of the
Springer resolution.  Its Borel--Moore homology carries the Hecke
algebra by convolution.  The Kazhdan--Lusztig basis arises from
intersection cohomology of orbit closures.  The representation
theory of $\mathfrak{S}$ is the representation theory of the Hecke
algebra: simple modules correspond to irreducible representations,
standard modules to Verma modules, the duality functor to the
Kazhdan--Lusztig involution.

\medskip

The bar complex of a chiral algebra $\cA$ lives on the product
$\FM_k(\C) \times \Conf_k(\R)$.  The differential extracts
collision residues on $\FM_k(\C)$; the coproduct deconcatenates
on $\Conf_k(\R)$.  The Koszul dual $\cA^!$ is the linear dual
of $\barB(\cA)$ (Verdier duality).  Line operators are
$\cA^!$-modules.  The representation theory of $\barB(\cA)$ is
the representation theory of the boundary: MC couplings are
line defects.

\medskip

\begin{center}
\small
\renewcommand{\arraystretch}{1.15}
\begin{tabular}{@{}lll@{}}
\textbf{Steinberg} & \textbf{Swiss-cheese} \\
\hline
$\widetilde\cN \to \cN$ (resolution)
  & $\FM_k(X) \to \Conf_k(X)$ (compactification) \\[3pt]
$\mathfrak{S} = \widetilde\cN
  \times_\cN \widetilde\cN$ (self-intersection)
  & $\barB(\cA) \otimes_{\barB(\cA\text{-bimod})} \barB(\cA)$
  (Hochschild) \\[3pt]
$H_*^{\mathrm{BM}}(\mathfrak{S})$ (Hecke algebra)
  & $\Zder(\cA)$ (derived center) \\[3pt]
KL basis (from IC)
  & PBW basis (from spectral sequence) \\[3pt]
Hecke modules
  & $\cA^!\text{-mod}$ (line operators) \\[3pt]
KL involution
  & Verdier/Koszul duality \\[3pt]
Cell structure
  & Shadow depth classification (G/L/C/M)
\end{tabular}
\end{center}

\medskip

The analogy becomes an identification at the level of
factorisation homology.  The factorisation homology of the bar
complex on a closed surface is the convolution product on the
self-intersection, which is the Hecke algebra.  The two
constructions produce the same algebra because both are computing
the endomorphism algebra of a compact generator in a factorisation
category on a configuration space.  The Fulton--MacPherson
compactification is the chiral analogue of the Springer resolution.

The cell structure of the Hecke algebra (Kazhdan--Lusztig cells,
Lusztig's $a$-function, the asymptotic Hecke algebra) should have
a shadow-tower analogue: the G/L/C/M classification of shadow
depth is the analogue of the cell decomposition, with shadow
depth playing the role of Lusztig's $a$-function.  Whether this
analogy extends to a theorem is open.


% =================================================================
\section{What remains}
\label{sec:what-remains}
% =================================================================

Everything in \S\ref{sec:two-directions}--\S\ref{sec:steinberg}
can be said in two sentences.

The bar complex of a chiral algebra on a curve carries a
differential (holomorphic direction) and a coproduct (topological
direction); together these form a Swiss-cheese algebra on
$\FM_k(\C) \times \Conf_k(\R)$, presenting the boundary theory
of a $3$d holomorphic-topological QFT whose bulk is the chiral
derived center, whose lines are modules for the Koszul dual, and
whose modularity arises from trace and clutching on the open
sector.  For gravity, the coproduct is primitive (ghost-number
obstruction from DS reduction), so all nonlinearity resides in the
$\Ainf$ products and the collision residue $r(z)$; for gauge
theory, the coproduct is nontrivial, and the full two-channel
structure is active.

That is what the $1$-form told us when we gave it a second
direction.  What follows is what it does when you push it further:
into the bordered moduli space, along the modular cooperad, past
the death of the derived category at genus $1$, through the
spectral Drinfeld strictification, into the gravitational
dichotomy at $c = 26$.  Each direction is a computation that
is either finished or honestly open.

\sep

The following structures are proved:
\begin{enumerate}[nosep]
\item The Swiss-cheese operad $\SCchtop$ and its
  homotopy-Koszulity.
\item The recognition theorem (Weiss cosheaf descent).
\item PVA descent D1--D6 (exchange cylinder, Stokes).
\item The four-stage open/closed architecture (Stages 1--3
  proved; Stage 4 programme).
\item The bulk-boundary-line triangle (boundary-linear sector
  proved; global conjectural).
\item Gravitational primitivity (principal DS; non-principal
  conjectural).
\item The benchmark: $V_k(\fsl_2)$ on $(\bP^1, \{0,\infty\})$,
  every component computed.
\item The annulus trace $\simeq \HH_*(\cC_{\mathrm{op}})$.
\item The completed platonic datum $\Pi^{\mathrm{oc}}_X(\cA)$.
\end{enumerate}

The following structures are open:
\begin{enumerate}[nosep]
\item The full modular cooperad on the trace.
\item The global bulk-boundary-line triangle outside the
  boundary-linear sector.
\item BV/BRST $=$ bar at genus $g \geq 1$.
\item Non-principal gravitational primitivity beyond hook-type.
\item The spectral Drinfeld strictification for Kac--Moody
  algebras with root multiplicities $> 1$.
\end{enumerate}

\sep

\begin{quote}\itshape
A $1$-form with a simple pole.

A relation among three.

A second direction: the coproduct.

The rest follows.
\end{quote}


% =================================================================
\section{The composite graviton}
\label{sec:composite-graviton}
% =================================================================

The Virasoro stress tensor
\[
T \;=\; \frac{1}{2(k+h^\vee)}\sum_a \normord{J^a J_a}
\]
is built from affine currents.  It is not fundamental.  The
graviton is composite.

The ghost sector of the BRST complex remembers this: the
coproduct is strictly primitive
(Theorem~\ref{thm:gravitational-primitivity-wn}; see the
gauge/gravity comparison in \S\ref{sec:coproduct}).

In gauge theory, the cross-term survives because $J^a$ is
fundamental (ghost number~$0$, no passage through~$h$):
\[
\Delta_z(T_{\mathrm{aff}}) \;=\;
\tau_z(T_{\mathrm{aff}}) \otimes 1
\;+\; 1 \otimes T_{\mathrm{aff}}
\;+\; \frac{1}{k+h^\vee}\,\Omega(w{+}z,\, w).
\]
The Casimir $\Omega = \sum_a J^a \otimes J_a$ is a production
vertex.  A gluon bound state dissociates into its constituents.

The $r$-matrix $r^{\mathrm{Vir}}(z) = (c/2)/z^3 + 2T/z$ carries
all gravitational scattering.  The cubic pole arises because
the bar kernel $d\log(z-w)$ absorbs one power from the quartic
OPE (see \S\ref{sec:graviton-composite} for explicit pole orders).


% =================================================================
\section{The holographic triangle}
\label{sec:holographic-triangle}
% =================================================================

Three vertices.

\medskip

\noindent
\textbf{Boundary}: $\cA$ (the chiral algebra).

\noindent
\textbf{Bulk}: $\Zder(\cA) = H^*(C^\bullet_{\mathrm{ch}}(\cA,\cA),\,\delta)$
(the chiral derived center).

\noindent
\textbf{Lines}: $\cA^!$-modules (on the chirally Koszul locus).

\sep

The bar complex $\barB(\cA)$ classifies couplings between boundary
and lines.  It is the universal twisting coalgebra: every map
$\alpha\colon \barB(\cA) \to \cA^!$ satisfying the twisting
morphism equation is a line operator.  On the Koszul locus:
\[
\MC\bigl(\Hom(\barB(\cA),\, \cA^!)\bigr)
\;\simeq\;
\Perf(\cA^!).
\]
This is the bar complex's role: a representability theorem.  It
classifies morphisms, not operators.

The derived center $\Zder(\cA)$ classifies bulk observables.
It acts on the boundary via the brace dg-algebra structure
on $C^\bullet_{\mathrm{ch}}(\cA,\cA)$, encoding how bulk fields
modify boundary correlators.  The
restriction $\Zder(\cA) \to \cA$ is well-defined.  The extension
$\cA \to \Zder(\cA)$ is not: bulk restricts to boundary, boundary
does not determine bulk.

\begin{warning}
$\barB(\cA) \neq \Zder(\cA)$.  The bar complex is a Fourier
kernel (Volume~I).  The derived center is a space of
endomorphisms.  One classifies morphisms; the other classifies
operators.  Their conflation is the most frequent source of false
statements in three-dimensional holomorphic-topological field
theory.
\end{warning}


% =================================================================
\section{Where modularity lives}
\label{sec:where-modularity-lives}
% =================================================================

As in \S\ref{sec:annulus}, modularity arises from trace and
clutching on the open sector, not as an axiom on the closed algebra.

\sep

Cut $S^1$ into two arcs.  Each arc is an interval $I \cong \R$.
Excision on the factorisation homology of $\cC_{\mathrm{op}}$:
\[
\int_{S^1}\!\cC_{\mathrm{op}}
\;\simeq\;
\cC_{\mathrm{op}}(I_1)
\;\otimes_{\cC_{\mathrm{op}}(I_1 \cap I_2)}^{\mathbb{L}}\;
\cC_{\mathrm{op}}(I_2).
\]
The two-arc cover produces the two-sided bar complex, which is the
cyclic bar complex.  Hochschild homology
$\HH_*(\cC_{\mathrm{op}})$ IS the annulus partition function.

The first modular shadow: $\kappa(\cA)$ appears as the leading
coefficient of the annulus trace.  The genus-$1$ curvature
$\dfib^{\,2} = \kappa \cdot \omega_1$ is the obstruction to
extending the genus-$0$ Swiss-cheese structure to the torus.  The
obstruction lives in the Hochschild homology of the open sector.
It does not live in any axiom on the closed algebra.

Clutching: the genus-$g$ surface $\Sigma_g$ is built from $3g-3$
pairs of pants, sewn along $3g-3$ circles.  Each sewing circle
contributes an $\HH_*$ factor.  The modular partition function is
the trace over all sewings:
\[
Z_g
\;=\;
\Tr_{\HH_*^{\otimes(3g-3)}}
\bigl(\text{clutching maps}\bigr).
\]
This is modularity from geometry, not from axioms.


% =================================================================
\section{The Koszul dual lives on the boundary}
\label{sec:koszul-dual-boundary}
% =================================================================

$D_{\Ran}(\barB(\cA)) \simeq \barB(\cA^!)$: the Verdier dual of the bar
coalgebra is the bar of the Koszul dual algebra~$\cA^!$
(Theorem~A intertwining).  The Koszul dual~$\cA^!$ lives on
the $\R$-direction (the boundary of the half-space
$\C_{\geq 0} \times \R$), not in the $\C \times \R$ bulk.

\sep

For Chern--Simons with gauge group $G$ at level~$k$:
$\cA = V_k(\fg)$, $\cA^! = V_{-k-2h^\vee}(\fg)$.  The
Feigin--Frenkel involution $k \mapsto -k - 2h^\vee$ is Verdier
duality on $\Ran(X)$.  At the critical level $k = -h^\vee$:
$V_{-h^\vee}(\fg)$ is commutative, bar cohomology is functions on
opers, and $\cA^!$ at the dual critical level is the
Feigin--Frenkel center.

For gravity: $\cA = \mathrm{Vir}_c$,
$\cA^! = \mathrm{Vir}_{26-c}$.

At $c = 26$: $\cA^! = \mathrm{Vir}_0$, uncurved
($\kappa = 0$).

At $c = 0$: $\cA = \mathrm{Vir}_0$, uncurved.

At $c = 13$: $\cA = \cA^!$.  Self-duality.

\medskip

The complementarity sum:
\[
\kappa(\mathrm{Vir}_c) + \kappa(\mathrm{Vir}_{26-c})
\;=\; \frac{c}{2} + \frac{26-c}{2} \;=\; 13 \;\neq\; 0.
\]
For Kac--Moody algebras: $\kappa + \kappa^! = 0$ by
Feigin--Frenkel (the involution $k \mapsto -k-2h^\vee$ is
anti-symmetric around $0$).  For Virasoro: the involution
$c \mapsto 26-c$ is anti-symmetric around $13$, not around $0$.
The midpoint matters.  It is 13, not 0.


% =================================================================
\section{The BTZ black hole from the $r$-matrix}
\label{sec:btz}
% =================================================================

The effective curvature of the matter-ghost system:
\[
\kappa_{\mathrm{eff}}
\;=\;
\kappa(\mathrm{Vir}_c) + \kappa_{\mathrm{ghost}}
\;=\;
\frac{c}{2} - 13
\;=\;
\frac{c-26}{2}.
\]
At $c = 26$: $\kappa_{\mathrm{eff}} = 0$.  No curvature.
The theory is topological.

At $c \neq 26$: the genus tower
$F_g = \kappa_{\mathrm{eff}} \cdot \lambda_g^{\mathrm{FP}}$.
The $\hat{A}$-genus generates all $F_g$ from $\kappa_{\mathrm{eff}}$:
\[
\sum_{g \geq 1} F_g\,\hbar^{2g}
\;=\;
\kappa_{\mathrm{eff}}
\bigl(\hat{A}(i\hbar) - 1\bigr).
\]
The Bernoulli numbers encode the higher-genus corrections.  Each
$F_g$ is positive (because $\hat{A}(ix) = (x/2)/\sin(x/2)$ has
all positive Taylor coefficients).

\sep

The collision residue
$r^{\mathrm{Vir}}(z) = (c/2)/z^3 + 2T/z$ controls BTZ
thermodynamics.  The leading coefficient $c/2 = \kappa$ is the
scalar that generates the Bekenstein--Hawking entropy at leading
order in $1/G$.  The subleading $2T/z$ is state-dependent braiding:
it distinguishes different BTZ states within the same theory.

The primitive coproduct is the algebraic statement that black holes
do not split.  The only channel is braiding.  The production
vertex $\Omega(w{+}z,w)$ of gauge theory (the Casimir cross-term)
is absent.  There is no splitting channel.

\begin{center}
\small
\renewcommand{\arraystretch}{1.15}
\begin{tabular}{@{}lcc@{}}
& $c \neq 26$ & $c = 26$ \\
\hline
$\kappa_{\mathrm{eff}}$ & $(c-26)/2 \neq 0$ & $0$ \\[3pt]
$F_1$ & $(c-26)/48$ & $0$ \\[3pt]
Genus tower & $\kappa_{\mathrm{eff}}\cdot(\hat{A}-1)$ & collapsed \\[3pt]
Regime & perturbative & topological
\end{tabular}
\end{center}


% =================================================================
\section{Five theories, one machine}
\label{sec:five-theories}
% =================================================================

\begin{center}
\small
\renewcommand{\arraystretch}{1.15}
\begin{tabular}{@{}lccccc@{}}
& CS & Gravity & Twisted hol.\ & M2 & M5 \\
\hline
$\cA$ & $V_k(\fg)$ & $\mathrm{Vir}_c$
  & $V_k(\fg)$ & DDCA$_K$ & $\cW_\infty(K)$ \\[3pt]
$\cA^!$ & $V_{-k-2h^\vee}(\fg)$ & $\mathrm{Vir}_{26-c}$
  & $V_{-k-2h^\vee}(\fg)$ & CE & ? \\[3pt]
$\kappa$ & $\dim(\fg)\frac{k+h^\vee}{2h^\vee}$
  & $c/2$ & $\dim(\fg)\frac{k+h^\vee}{2h^\vee}$
  & $K(K{+}1)/2$ & $\varrho_\infty c$ \\[3pt]
Class & L & M & L & L & M \\[3pt]
$\Delta_z$ & nontrivial & primitive & nontrivial
  & nontrivial & primitive \\[3pt]
$\kappa + \kappa^!$ & $0$ & $13$ & $0$ & $0$
  & $\varrho K_N$
\end{tabular}
\end{center}

\medskip

The machine does not know which theory it is computing.  It
receives: a chiral algebra, its OPE, and the bar construction.
The five theories are five inputs to the same functor.

Chern--Simons and twisted holography are class~L: the $r$-matrix
carries all structure, the coproduct is nontrivial, the shadow
tower terminates at arity~$3$.  Gravity and M5 are class~M:
infinite shadow depth, primitive coproduct,
$Q^{\mathrm{ct}} \neq 0$.

\sep

The M2--M5 duality: the M2~algebra (DDCA, class~L) and the
M5~algebra ($\cW_\infty$, class~M) are related by the Miura
transformation.  DS~reduction converts class~L to class~M.
The reduction IS the depth increase.  The ghost-number filtration
that kills the coproduct of M5 is the same filtration that was
absent for M2.


% =================================================================
\section{The four-stage architecture, revisited}
\label{sec:four-stages-revisited}
% =================================================================

Stage~1 (local one-colour): define the Swiss-cheese algebra.
Prove the universal acting-bulk theorem
$U(\cA) = (C^*_{\mathrm{ch}}(\cA,\cA),\, \cA)$.
Everything succeeds or fails here.

\medskip

Stage~2 (open primitive): the primitive object is
$\cC_{\mathrm{op}}$, not $A_b$.  The boundary algebra
$A_b = \End(b)$ is a chart of a compact generator $b$ in
$\cC_{\mathrm{op}}$.  Different generators produce
Morita-equivalent charts.  The category is the invariant.

\medskip

Stage~3 (globalise): move from $\C \times \R$ to tangential log
curves $(X, D, \tau)$.  The real oriented blowup
$\pi\colon \widetilde{X}_D \to X$ replaces each puncture $p_i$ by
a circle $S^1_{p_i}$; the tangential basepoint $\tau_{p_i}$ cuts
this circle into an interval $I_{p_i} \cong \R$.  Open insertions
live on the interval.  Bulk insertions live in the interior.  The
global theory is a factorisation algebra on a bordered Riemann
surface.

\medskip

Stage~4 (modularise): introduce $\Theta^{\mathrm{oc}}$, traces,
annulus/Hochschild chains, clutching, quartic resonance classes.
Only after stages~1--3 are solid.

\sep

Each stage depends on its predecessor.  Discussion of stage~$k$
before establishment of stage~$k-1$ is not premature; it is wrong.


% =================================================================
\section{What the bar complex knows about quantum gravity}
\label{sec:bar-knows-gravity}
% =================================================================

The bar complex $\barB(\mathrm{Vir}_c)$ is a factorisation
coalgebra on $\Ran(X)$.  Its differential $d_{\barB}$ extracts
collision residues of the Virasoro OPE.  Its coproduct~$\Delta$
deconcatenates ordered tuples along $\R$.

At genus~$0$: $d_{\barB}^2 = 0$ by Arnold.  The bar complex is
an honest chain complex.  Its cohomology has dimensions
\[
1,\; 2,\; 5,\; 12,\; 30,\; 76,\; 196,\; 512,\; 1353,\; \ldots
\]
The generating function satisfies
$x^3 P^2 - (1-2x)(1-x^2)P + (1-x)(1-x^2) = 0$.
Discriminant: $(1-x)^2(1-3x)(1+x)$.  Growth rate~$3$.

\sep

At genus~$1$:
$\dfib^{\,2} = (c/2) \cdot \omega_1 \neq 0$.
The bar complex is curved.  The curvature $c/2$ is the modular
characteristic.

At genus~$g$:
\[
F_g
\;=\;
\frac{c}{2} \cdot \lambda_g^{\mathrm{FP}}
\;=\;
\frac{c}{2} \cdot
\frac{(2^{2g-1} - 1)\,|B_{2g}|}{(2g)!\cdot 2^{2g-1}}.
\]
The Bernoulli numbers.  The $\hat{A}$-genus.  From index theory.
Nothing in the Virasoro OPE predicts this.  It lives in the
geometry of $\overline\cM_g$.

\sep

The shadow tower
\begin{align*}
S_2 &= c/2, \\
S_3 &= 2, \\
S_4 &= 10\big/\bigl[c(5c+22)\bigr], \\
S_5 &= -48\big/\bigl[c^2(5c+22)\bigr],
\end{align*}
and all higher $S_r$, are the Taylor coefficients of $\sqrt{Q}$
where
\[
Q_{\mathrm{Vir}}(t)
\;=\;
c^2 + 12ct
+ \frac{180c + 872}{5c + 22}\,t^2.
\]
One quadratic.  One square root.  An algebraic function of
degree~$2$.

\medskip

Three numbers determine the entire infinite tower of higher-genus
corrections to quantum gravity: $\kappa = c/2$, $\alpha = 2$ (the
cubic shadow), and $S_4 = 10/[c(5c+22)]$ (the quartic contact
invariant).  The shadow metric $Q$ is constructed from these three.
Its discriminant $\Delta = 8\kappa S_4$ controls the convergence
radius of the shadow tower.  When $\Delta \neq 0$ (all
$c \neq 0$): the tower is infinite, class~M.  The bar complex
knows the exact answer at every genus, as the coefficients of an
algebraic function.


% =================================================================
\section{The arithmetic structure of quantum gravity}
\label{sec:arithmetic-gravity}
% =================================================================

The gravitational shadow tower is the arithmetic of the quadratic
extension $K = \Q(c)(t)(\sqrt{Q_{\mathrm{Vir}}})$ over the
function field $F = \Q(c)(t)$.  The shadow metric
\[
  Q_{\mathrm{Vir}}(t)
  \;=\;
  c^2 + 12ct + \frac{180c + 872}{5c + 22}\,t^2
\]
is a quadratic polynomial whose square root generates the
infinite tower of higher-genus corrections to three-dimensional
quantum gravity.

\sep

\noindent\textbf{The discriminant.}\enspace
$\mathrm{disc}(Q_{\mathrm{Vir}}) = -320\,c^2/(5c{+}22)$.
Negative for all $c > 0$: the roots of $Q_{\mathrm{Vir}}$
are complex conjugate, the extension $K/F$ is genuine, and
the shadow tower is infinite (class~M).

The discriminant controls three quantities at once:
\begin{enumerate}[label=(\roman*),nosep]
\item The depth of the gravitational perturbation series
  (infinite for all $c \neq 0$, because $\Delta = 40/(5c{+}22)
  \neq 0$).
\item The convergence radius of the shadow tower: the growth
  rate $\rho = \sqrt{36 + 80/(5c{+}22)}\,/\,c$ is the reciprocal
  of the distance from the origin to the nearest branch point
  of $\sqrt{Q_{\mathrm{Vir}}}$.
\item The complementarity with the dual theory at $c' = 26 - c$:
  $\Delta(c) + \Delta(26{-}c) = 6960/[(5c{+}22)(152{-}5c)]$,
  with universal numerator.
\end{enumerate}

\sep

\noindent\textbf{The gravitational number field.}\enspace
The Bekenstein--Hawking entropy $S_{\mathrm{BH}} = 2\pi\sqrt{ch/6}$
derives from $\kappa = c/2$, the leading coefficient of
$Q_{\mathrm{Vir}}$.  This coefficient is the constant term
(up to a square): $Q_{\mathrm{Vir}}(0) = c^2 = 4\kappa^2$.
The leading OPE datum seeds the quadratic extension.

The growth rate $\rho$ measures how fast the higher-genus
corrections to the gravitational partition function grow.
When $\rho < 1$, the shadow tower converges: the shadow
partition function $Z^{\mathrm{sh}} = \sum Z_g^{\mathrm{sh}}
\hbar^{2g}$ exists as a convergent series.  This convergence
is the ``splitting'' of the gravitational prime at central
charge~$c$: the quadratic extension has a rational point in
a neighbourhood of the origin.

Setting $\rho = 1$ yields the critical cubic
$5c^3 + 22c^2 - 180c - 872 = 0$ with positive root
$c^* \approx 6.125$.  For $c > c^*$: the shadow tower
converges ($\rho < 1$, the gravitational perturbation
series sums).  For $0 < c < c^*$: the shadow tower diverges
($\rho > 1$, the perturbation series requires resummation).
The critical cubic is the Minkowski bound of the gravitational
number field: it separates the convergent regime from the
divergent regime.

\sep

\noindent\textbf{The three critical values.}\enspace
\begin{center}
\small
\renewcommand{\arraystretch}{1.15}
\begin{tabular}{@{}ccl@{}}
$c$ & \textbf{What happens} & \textbf{Arithmetic meaning} \\
\hline
$0$ & $\kappa = 0$, $Q$ degenerates & ramified prime
  (shadow tower is undefined) \\[3pt]
$13$ & $\cA = \cA^!$, self-dual
  & Galois involution is trivial on coefficients;
  $K = K'$ \\[3pt]
$26$ & $\kappa_{\mathrm{eff}} = 0$, topological & dual
  extension $K'$ degenerates ($\kappa' = 0$)
\end{tabular}
\end{center}

\medskip

At $c = 0$: the constant term of $Q_{\mathrm{Vir}}$ vanishes.
The quadratic polynomial has a double root at $t = 0$.  The
shadow connection $\nabla^{\mathrm{sh}}$ has an irregular
singularity.  The extension ramifies.

At $c = 13$: the theory is self-dual ($\mathrm{Vir}_{13}^!
= \mathrm{Vir}_{13}$).  The Koszul involution $c \mapsto 26-c$
fixes $c = 13$.  The two quadratic extensions $K$ and $K'$
coincide.  The complementarity sum $\kappa + \kappa' = 13$
does not vanish: self-duality is a symmetry enhancement, not
anomaly cancellation.  The growth rate $\rho(13) \approx 0.467$:
the shadow tower converges at the self-dual point.

At $c = 26$: the effective curvature $\kappa_{\mathrm{eff}}
= (c{-}26)/2 = 0$ vanishes.  The dual algebra
$\mathrm{Vir}_0$ has $\kappa' = 0$: its shadow extension
ramifies.  The gravitational theory becomes topological.  The
Clifford algebra degenerates from a matrix algebra
($2^g$-dimensional) to an exterior algebra
($2^{2g}$-dimensional).  The BTZ partition function collapses:
$Z_{\mathrm{grav}} = 1$.

These three values are not phases of one theory.  They are three
qualitatively distinct points in the arithmetic landscape of the
quadratic extension: a ramified prime ($c = 0$), a self-dual
fixed point ($c = 13$), and a degeneration of the dual
($c = 26$).  They are separated by $13$ units of central charge.

\sep

\noindent\textbf{The ramified primes.}\enspace
The denominator theorem (Volume~I) gives
$\mathrm{denom}(S_r)
= (\mathrm{const})\cdot c^{r-3}(5c{+}22)^{\lfloor(r-2)/2\rfloor}$.
Two primes ramify:
\begin{enumerate}[label=(\roman*),nosep]
\item $c = 0$: the leading OPE coefficient vanishes.  No
  Bekenstein--Hawking entropy.  No bar curvature.  The shadow
  metric $Q(0,t) = 872\,t^2/(5 \cdot 22) = 872\,t^2/110$
  degenerates.
\item $c = -22/5$: the weight-$4$ Gram determinant
  $\det G_4 = c^2(5c{+}22)/2$ vanishes.  A null vector at
  weight~$4$ collapses the quasi-primary $\Lambda$ onto
  $\partial^2 T$.  The quartic contact invariant
  $Q^{\mathrm{ct}} = 10/[c(5c{+}22)]$ has a pole.
\end{enumerate}
Every other value of~$c$ is unramified: the shadow
coefficients are finite, the quadratic extension is
well-defined, and the tower runs.

No higher Kac determinant appears.  The gravitational number
field has exactly two ramified primes, both visible from the
weight-$4$ representation theory.  The infinite tower of
higher-genus gravitational corrections is controlled by a
single $2 \times 2$ Gram matrix.

\sep

\noindent\textbf{The unit group and the BTZ partition function.}\enspace
The growth rate
$\rho = \sqrt{36 + 80/(5c{+}22)}\,/\,c$ is the fundamental
unit of the gravitational number field.  It controls the BTZ
partition function through the radius of convergence of the
shadow tower:
\begin{itemize}[nosep]
\item $\rho < 1$ ($c > c^*$): the shadow tower converges.  The
  gravitational path integral exists non-perturbatively as a
  convergent shadow series.  The shadow partition function
  $Z^{\mathrm{sh}}$ is the value of a convergent double series
  (genus and arity).
\item $\rho > 1$ ($0 < c < c^*$): the shadow tower diverges.
  The shadow coefficients grow exponentially.  Resummation
  (Borel transform of $\sqrt{Q_{\mathrm{Vir}}}$) is required.
  The Stokes constants of the shadow tower are the
  non-perturbative content.
\item $\rho = 1$ ($c = c^*$): the critical cubic.  The branch
  point of $\sqrt{Q_{\mathrm{Vir}}}$ lies exactly on the unit
  circle.  The shadow tower is marginally divergent.
\end{itemize}

\begin{center}
\small
\renewcommand{\arraystretch}{1.15}
\begin{tabular}{@{}rrl@{}}
$c$ & $\rho$ & \textbf{Regime} \\
\hline
$1$ & $2.99$ & divergent (shadow) \\
$c^* \approx 6.125$ & $1.00$ & critical \\
$13$ & $0.467$ & convergent \\
$26$ & $0.234$ & convergent (dual topological) \\
$100$ & $0.061$ & strongly convergent
\end{tabular}
\end{center}

\medskip

The string partition function $Z_{\mathrm{string}}$ (genus
factorial growth $(2g)!$) always diverges.  The shadow partition
function $Z^{\mathrm{sh}}$ (Bernoulli decay $1/(2\pi)^{2g}$
times shadow growth $\rho^r r^{-5/2}$) converges for $c > c^*$.
The shadow tower is better-behaved than the string because it
sums over arities within a fixed genus, while the string sums
over genera with unrestricted growth.

\sep

\noindent\textbf{The Galois involution.}\enspace
The monodromy $-1$ of the shadow connection
$\nabla^{\mathrm{sh}} = d - Q'/2Q\,dt$ around each branch
point of $\sqrt{Q_{\mathrm{Vir}}}$ is the nontrivial element
of $\mathrm{Gal}(K/F)$.  In the gravitational context: the
Koszul sign is the Galois involution of the number field.  The
sign alternation in the shadow tower $S_r$ (positive for even
$r$, approximately negative for odd $r$, with beat period
$\pi/(\pi - \theta)$ where $\theta = \arg(t_+)$) is the
manifestation of the Galois action on the Taylor coefficients
of $\sqrt{Q_{\mathrm{Vir}}}$.

This is the deepest structural content of the bar
construction's desuspension: the sign $(-1)^n$ at bar degree
$n$ is the Galois automorphism of a quadratic extension,
evaluated at the $n$-th coefficient.

\sep

The gravitational shadow tower is not like a quadratic
extension.  It is one, over the function field $\Q(c)(t)$
instead of $\Q$.  The Bekenstein--Hawking entropy is the leading
coefficient.  The convergence of the perturbation series is the
splitting behaviour.  The Koszul sign is the Galois involution.
The critical cubic is the Minkowski bound.  The two ramified
primes are the Kac determinant zeros.  The complementarity
identity is the duality constraint on discriminants.  The
gravitational partition function is the norm form of the
extension, evaluated at a specific point.

\begin{quote}\itshape
A $1$-form with a simple pole.

A relation among three.

A second direction: the coproduct.

A quadratic extension: the shadow.

The rest follows.
\end{quote}


% =================================================================
\section{Gravity at rational central charge}
\label{sec:gravity-rational-c}
% =================================================================

The generic gravitational number field
$K = \Q(c)(t)(\sqrt{Q_{\mathrm{Vir}}})$ of the preceding
section lives over the function field $\Q(c)$.  When the
central charge specialises to a rational number $c_0$, the
field specialises to a concrete quadratic extension of~$\Q$.

For a minimal model $c_{p,q} = 1 - 6(p-q)^2/(pq)$,
the critical discriminant
$\Delta(c_{p,q}) = 40/(5c_{p,q} + 22)$ is a specific
rational number, and
$K(c_{p,q}) := \Q(\sqrt{\Delta(c_{p,q})})$
is a specific real quadratic number field.  Gravity at
each minimal model selects a number field.

\sep

\noindent\textbf{The physical picture.}\enspace
The Bekenstein--Hawking entropy $\kappa = c/2$ at a rational
central charge $c_{p,q}$ is a rational number.  The
higher-genus corrections $F_g = \kappa \cdot \lambda_g^{\mathrm{FP}}$
are rational multiples of~$\kappa$, hence rational.  The
ENTIRE perturbative gravitational partition function
$\sum F_g \hbar^{2g}$ has rational coefficients.

But the shadow tower is different.  The Taylor coefficients
$S_r$ of $\sqrt{Q_L}$ are rational (they satisfy a rational
recursion), yet the generating function $H(t) = t^2\sqrt{Q_L(t)}$
is genuinely irrational: it involves $\sqrt{\Delta}$, which
lies outside $\Q$ for all $c > 0$.  The irrationality is the
arithmetic content of the shadow tower.  The bar complex of
the Virasoro algebra, specialised to a minimal model, produces
an algebraic number that no finite truncation of the
perturbative expansion can detect.

This is not a deficiency of perturbation theory.  The function
$H(t)$ is algebraic of degree~$2$; its Taylor coefficients are
all individually rational.  The irrationality is non-perturbative
in the arity variable $t$, not in the genus variable $\hbar$.
It is visible in the branch structure of $\sqrt{Q_L}$, not in
any single coefficient.

\sep

\noindent\textbf{Computation.}\enspace
Write $5c_0 + 22 = a/b$ in lowest terms.  Then
$\Delta(c_0) = 40b/a$, and the specialised field is
$K(c_0) = \Q(\sqrt{\mathrm{sqfree}(40ab)})$.

\medskip

\noindent\textit{$c = 1/2$ (Ising model).}\enspace
$5c + 22 = 49/2$.  $\Delta = 80/49$.
$\sqrt{80/49} = 4\sqrt{5}/7$.
The Ising gravity lives over $\Q(\sqrt{5})$.

The golden ratio $(1+\sqrt{5})/2$ is the fundamental unit
of this field, with norm $-1$.  The class number is $h = 1$:
the ring of integers $\Z[(1+\sqrt{5})/2]$ is a principal ideal
domain.  The Ising model produces the simplest possible
arithmetic.  One is tempted to say that the Ising model is
``arithmetically trivial,'' as it is physically the simplest
unitary minimal model.

The shadow growth rate is $\rho(1/2) \approx 12.53$: the
shadow tower diverges violently ($c = 1/2$ is far below
the critical cubic $c^* \approx 6.125$).  The Ising gravity
requires resummation.

\medskip

\noindent\textit{$c = 7/10$ (tricritical Ising).}\enspace
$5c + 22 = 51/2$.  $\Delta = 80/51$.
$K(7/10) = \Q(\sqrt{255})$ where $255 = 3 \cdot 5 \cdot 17$.
Discriminant $1020$.  Class number $h = 4$: a genuine failure
of unique factorisation.

\medskip

\noindent\textit{$c = 4/5$ (three-state Potts).}\enspace
$K(4/5) = \Q(\sqrt{65})$.  Discriminant~$65$.  $h = 2$.
The Potts-model gravity has a non-trivial ideal class group
$\mathrm{Cl}(K) \cong \Z/2$.

\medskip

\noindent\textit{$c = 13$ (self-dual).}\enspace
$K(13) = \Q(\sqrt{870})$ where $870 = 2 \cdot 3 \cdot 5 \cdot 29$.
Discriminant $3480$.  $h = 8$.

The self-dual gravity has the largest class number of any
integer central charge in $[0, 26]$.  It sees the prime~$29$,
inherited from the complementarity constant
$6960 = 2^4 \cdot 3 \cdot 5 \cdot 29$.  No other central
charge in the standard range detects this prime.

\medskip

\noindent\textit{$c = 26$ (critical string).}\enspace
$K(26) = \Q(\sqrt{95})$.  Discriminant $380$.  $h = 2$.
$\rho(26) \approx 0.232$: convergent.  This is the
gravitational number field of the bosonic string at the
critical dimension.

\begin{center}
\small
\renewcommand{\arraystretch}{1.15}
\begin{tabular}{@{}lccrrcr@{}}
$c$
  & $\Delta(c)$
  & $d(c)$
  & $\mathrm{disc}$
  & $h$
  & $\varepsilon$
  & $\rho(c)$ \\
\hline
$1/2$ & $80/49$ & $5$ & $5$ & $1$
  & $(1{+}\sqrt{5})/2$ & $12.53$ \\[2pt]
$7/10$ & $80/51$ & $255$ & $1020$ & $4$
  & $16{+}\sqrt{255}$ & $8.94$ \\[2pt]
$4/5$ & $20/13$ & $65$ & $65$ & $2$
  & $8{+}\sqrt{65}$ & $7.81$ \\[2pt]
$1$ & $40/27$ & $30$ & $120$ & $2$
  & $11{+}2\sqrt{30}$ & $6.24$ \\[2pt]
$13$ & $40/87$ & $870$ & $3480$ & $8$
  & $59{+}2\sqrt{870}$ & $0.47$ \\[2pt]
$25$ & $40/147$ & $30$ & $120$ & $2$
  & $11{+}2\sqrt{30}$ & $0.24$ \\[2pt]
$26$ & $5/19$ & $95$ & $380$ & $2$
  & $39{+}4\sqrt{95}$ & $0.23$
\end{tabular}
\end{center}

\sep

\noindent\textbf{Koszul duality and the field pair.}\enspace
The Koszul pair $(c, 26-c)$ produces two fields $K(c)$
and $K(26-c)$.  These coincide iff
$(5c+22)/(152-5c) \in \Q^{\times 2}$.

At $c = 1$ and $c = 25$: the ratio is $27/147 = (3/7)^2$.
Both give $\Q(\sqrt{30})$.

At $c = 13$: the ratio is $87/87 = 1$.  Self-duality forces
field-preservation.

For the Ising model: the ratio is $49/299 = 49/(13 \cdot 23)$,
which is not a perfect square.  The Ising gravity at $c = 1/2$
lives over $\Q(\sqrt{5})$; its Koszul dual at $c = 51/2$
lives over $\Q(\sqrt{1495})$ where $1495 = 5 \cdot 13 \cdot 23$.
Different fields, different class numbers, different arithmetic.
Koszul duality does not preserve the gravitational number field
at the Ising point.

This failure is generic.  Among all minimal models $M(p,q)$
with $p < 10$, none has the Koszul pair giving the same field.
The criterion $(5c+22)/(152-5c) \in \Q^{\times 2}$ requires
an arithmetic coincidence that the minimal model spectrum
does not provide.

\sep

\noindent\textbf{What the class number measures.}\enspace
The class number $h$ of $K(c)$ is a measure of the
arithmetic complexity of gravity at central charge~$c$.
In algebraic number theory, $h$ counts the failure of
unique factorisation in the ring of integers $\cO_K$.

In the gravitational context: $h = 1$ means the shadow tower's
arithmetic is ``principal,'' governed by a single generator.
$h > 1$ means the shadow arithmetic requires multiple ideal
classes.  Whether this has a direct physical interpretation
(a counting of distinct gravitational saddle classes, perhaps)
is not known.

\begin{center}
\small
\renewcommand{\arraystretch}{1.15}
\begin{tabular}{@{}lcl@{}}
$c$ & $h$ & \textbf{Interpretation} \\
\hline
$1/2$ (Ising) & $1$ & principal: simplest arithmetic \\[2pt]
$4/5$ (Potts) & $2$ & one non-trivial ideal class \\[2pt]
$7/10$ (tricrit.) & $4$ & four ideal classes \\[2pt]
$13$ (self-dual) & $8$ & most complex in standard range
\end{tabular}
\end{center}

\sep

\noindent\textbf{What fails.}\enspace
The growth rate $\rho(c)$ is a smooth algebraic function
of~$c$.  The fundamental unit $\varepsilon(d(c))$ of $K(c)$
depends on $c$ through the squarefree part $d(c)$, which is a
wildly discontinuous function of the central charge.  No
algebraic relation connects $\rho$ and $\varepsilon$.

A conjecture that the fundamental unit of the gravitational
number field controls the shadow growth rate was tested by
computing both quantities at $c = 1/2, 7/10, 4/5, 1, 13, 25, 26$
and found to be false.  At $c = 1/2$: $\rho \approx 12.53$
while $\varepsilon = (1+\sqrt{5})/2 \approx 1.618$.  At
$c = 13$: $\rho \approx 0.467$ while
$\varepsilon = 59 + 2\sqrt{870} \approx 117.99$.  The two
quantities vary in opposite directions.

The failure is structural: $\rho(c)$ is a smooth function
of~$c$, while $\varepsilon(d(c))$ is a step function.
A smooth function cannot equal a step function.

\sep

\noindent\textbf{What remains.}\enspace
The specialisation $c \mapsto K(c)$ is a new functor from
(a subset of) the rational points of the Virasoro central
charge line to the set of real quadratic number fields.  It
sends the OPE data at genus~$0$ (the stress tensor and its
weight-$4$ descendant) to a number field whose arithmetic
is entirely determined by the bar complex.

The family $\{K(c_{p,q})\}_{(p,q)}$ attached to the minimal
model tower constitutes a new arithmetic invariant of 2d~CFT.
Whether this family has structure beyond what is visible in
the table above is an open question.

The first instance of the specialisation functor applied to
gravity is the Ising model.  The shadow tower of the
Ising model lives over $\Q(\sqrt{5})$, the field of the
golden ratio.  This is the number field with class number~$1$
and the smallest possible discriminant among real quadratic
fields with fundamental unit of norm $-1$.  The simplest
unitary minimal model produces the simplest real quadratic
field.

Whether this simplicity persists along the minimal model tower,
or whether more complex arithmetic emerges, is a question at
the intersection of vertex algebra theory and algebraic number
theory that the bar construction makes precise.

\begin{quote}\itshape
The bar complex is a 1-form.

Its square root is a number field.

Gravity at rational central charge is arithmetic.
\end{quote}


% =================================================================
\section{Fourier duality from operads to $L$-functions}
\label{sec:fourier-operads-L}
% =================================================================

The bar construction is a Fourier transform
(\S\ref{sec:fourier}, Volume~I).  The shadow tower is a
quadratic extension (\S\ref{sec:arithmetic-gravity}).  The
Dirichlet-sewing lift $S_\cA(u)$ is a product of zeta functions
(\S\ref{sec:prime-factorisation}, Volume~I).  These three objects
arise from the same chiral algebra.  This section asks: do they
form a ``descent chain'' in which a single Fourier duality
descends from operadic algebra to analytic number theory?

The answer is: partially.  The chain has five levels and three
distinct involutions.  The first two links are theorems.  The
third is broken.  The fourth is lattice-specific.  The claim
``they are the same involution'' is false.

\subsection{The three involutions}

The physical content of the bar construction is that it extracts
the OPE residues of a chiral algebra into a nilpotent differential,
converting the multiplicative structure (operator product) into an
additive structure (differential).  This extraction carries three
symmetries:

\begin{enumerate}[label=\textup{(\roman*)}]
\item \textbf{Verdier duality} $\bD_{\Ran}$ on $\Ran(X)$.
  Exchanges $\barB(\cA)$ and $\barB(\cA^!)$ (Theorem~A,
  Volume~I).  This is the Fourier involution: it intertwines
  convolution with pointwise operations on factorisation
  (co)algebras.  Its kernel is $\eta_{ij} = d\log(z_i - z_j)$.

  Fixed locus: self-dual factorisation coalgebras.

\item \textbf{Galois involution} $\sigma$ of the shadow
  extension $K_L = \Q(c)(t)(\sqrt{Q_L})$.  Sends
  $\sqrt{Q_L} \mapsto -\sqrt{Q_L}$.  This is the monodromy
  of the shadow connection $\nabla^{\mathrm{sh}}$: transport
  around a branch point of $\sqrt{Q_L}$ returns $-\sqrt{Q_L}$.
  It is the Koszul sign $(-1)^n$ at bar degree~$n$.

  Fixed locus: the base field $F = \Q(c)(t)$.

\item \textbf{Koszul involution} $\cA \mapsto \cA^!$ on
  families of chiral algebras.  For Virasoro: $c \mapsto 26 - c$.
  For Kac--Moody: $k \mapsto -k - 2h^\vee$.  This acts on the
  \emph{coefficients} of the shadow metric $Q_L$, not on
  $\sqrt{Q_L}$ itself.

  Fixed locus: self-dual algebras ($c = 13$ for Virasoro).
\end{enumerate}

These three involutions are related but distinct.  Verdier
duality produces the Koszul involution as its shadow on moduli:
$\bD_{\Ran}(\barB(\cA)) \simeq \barB(\cA^!)$ implies that
$\cA \mapsto \cA^!$ is the trace of Verdier duality on isomorphism
classes.  The Galois involution is intrinsic to a \emph{single}
algebra's shadow tower and has no direct relation to Verdier
duality.  It arises because the MC equation
$D\Theta + \tfrac{1}{2}[\Theta, \Theta] = 0$, projected to the
scalar arity line, reduces to a Riccati equation whose solution
$H(t) = t^2\sqrt{Q_L}$ is algebraic of degree~$2$---not higher,
because the Riccati is quadratic.

The claim ``they are the same involution (Fourier) acting at
different levels'' conflates (i)--(iii).  The correct statement:
they are three $\Z/2$-actions on three different categories, all
originating from the bar construction but not related by any
``descent.''  At the self-dual point $c = 13$, all three have
compatible fixed-point behaviour (the Koszul pair coincides, the
shadow extensions $K_L$ and $K_{L!}$ agree, and the Verdier dual
coalgebra is isomorphic to the original).  At generic $c$, they
diverge: the Galois involution of $K_L$ has nothing to do with
the Koszul involution on the $c$-line.

\subsection{Where the chain breaks}

The descent chain has five levels:

\begin{center}
\small
\renewcommand{\arraystretch}{1.15}
\begin{tabular}{@{}cll@{}}
\textbf{Level} & \textbf{Object} & \textbf{Status} \\
\hline
0 & Koszul duality $\cA \mapsto \cA^!$
  & Proved (Theorem~A) \\
1 & Bar as Fourier: $\barB_X$, inversion, Plancherel
  & Proved (Theorems A,B,C) \\
2 & Shadow tower as quadratic extension
  & Proved \\
3 & Dirichlet-sewing lift $S_\cA(u)$
  & Defined, but independent of Level 2 \\
4 & Zeros of $\varepsilon^c_s$ on critical lines
  & Lattice only; structural obstruction
\end{tabular}
\end{center}

The break occurs between Levels 2 and 3.  The shadow tower
is determined by $(\kappa, \alpha, S_4)$---the arity-$2$, $3$,
$4$ OPE data---and lives over $\Q(c)(t)$.  The Dirichlet-sewing
lift is determined by the weight multiset $W(\cA) = \{w_i\}$
and lives over the complex $u$-plane.  These are independent
invariants: two algebras with the same weight multiset but
different OPE structure constants have the same $S_\cA(u)$ but
different shadow towers.  Conversely, two algebras with the same
$(\kappa, \alpha, S_4)$ but different weight multiplicities have
the same shadow tower but different $S_\cA(u)$.

\subsection{The structural obstruction}

For lattice VOAs, the constrained Epstein zeta $\varepsilon^c_s$
factors into $L$-functions because the theta function $\Theta_\Lambda$
decomposes under the Hecke algebra into Eisenstein series and
eigenforms.  This decomposition uses three properties special to
the lattice setting: the theta function is a classical modular form,
the Hecke decomposition is a finite sum, and the Fourier coefficients
are lattice representation numbers with multiplicative structure.

None of these holds for the Virasoro algebra.  The Virasoro
character $q^{-c/24}\prod(1-q^n)^{-1}$ is not a modular form of
finite weight.  The shadow tower still runs (class~M, infinite
depth), but the arithmetic sieve does not apply.  The MC equation
constrains the spectral coefficients on the \emph{real} spectral
axis; the zeta zeros live at \emph{complex} spectral parameters.
This separation is structural
(Remark~\ref{rem:structural-obstruction}, Volume~I): algebraic
constraints on the spectral line cannot reach the scattering poles
without analytic continuation of the spectral data off the real
axis.

Concretely: the modified Li coefficients
$\tilde\lambda_n(\cH)$ for the Heisenberg sewing lift
$S_\cH(u) = \zeta(u)\zeta(u{+}1)$ are positive for
$n = 1, \ldots, 6$ and \textbf{negative at $n = 7$}
(\texttt{compute/tests/test\_euler\_koszul\_engine.py}, line~394).
Positivity of all Li coefficients would be equivalent to GRH for
$\zeta(u)\zeta(u{+}1)$.  The sign change at $n = 7$ shows that the
sewing-lift Li coefficients do not satisfy a ``shadow RH.''  (This
does not contradict RH for $\zeta$ alone; the sign change is an
artefact of the product and the regularisation.)

\subsection{What the bar construction does provide}

The descent chain, as a linear sequence, is broken.  But the
bar construction does provide something more interesting: a
\emph{fan} of three independent projections from a single
MC element.

\medskip

\begin{center}
\small
\begin{tabular}{@{}lll@{}}
\textbf{Projection} & \textbf{Output} & \textbf{Content} \\
\hline
Genus tower
  & $F_g = \kappa \cdot \lambda_g^{\mathrm{FP}}$
  & Topological (Mumford classes) \\[3pt]
Shadow tower
  & $\{S_r\}_{r \geq 2} = [t^{r-2}]\sqrt{Q_L}$
  & Arithmetic (quadratic extension, discriminant) \\[3pt]
Koszul dual
  & $\cA^! = (\barB(\cA))^\vee$
  & Algebraic (Verdier dual)
\end{tabular}
\end{center}

\medskip

These three projections are facets of $\Theta_\cA
\in \mathrm{MC}(\gAmod)$, the universal MC element in the
modular cyclic deformation complex.  The genus tower is the
scalar projection of $\Theta_\cA$ to each genus.  The shadow
tower is the arity projection of $\Theta_\cA$ on the primary
line.  The Koszul dual is the genus-$0$, arity-$2$ projection
(the binary pairing that determines $\cA^!$).

In the gravitational context ($\cA = \mathrm{Vir}_c$):
\begin{itemize}[nosep]
\item The genus tower gives $F_g = (c/2)\lambda_g^{\mathrm{FP}}$:
  the Bernoulli numbers.  Higher-genus gravitational partition
  functions.
\item The shadow tower gives $\sqrt{Q_{\mathrm{Vir}}}$: the
  quadratic extension $K = \Q(c)(t)(\sqrt{Q})$ with discriminant
  $-320c^2/(5c{+}22)$.  The arithmetic of quantum gravity.
\item The Koszul dual gives $\mathrm{Vir}_{26-c}$: the dual
  gravitational theory.  $c + c' = 26$ is Freudenthal--de~Vries.
\end{itemize}

All three are proved.  All three are outputs of the single
object $\Theta_{\mathrm{Vir}}$.  But they are three outputs, not
five levels of one chain.

\subsection{The correct picture}

The bar construction is a Fourier transform.  This is a theorem,
not a metaphor (Volume~I, Theorems~A, B, C).  The shadow tower
is a quadratic extension.  This is a theorem, not an analogy
(\S\ref{sec:arithmetic-gravity}).  But the claim that these two
facts are levels of a ``descent chain'' that continues to
$L$-functions and the Riemann hypothesis is false.

The reason it is false is not a failure of imagination; it is a
theorem.  The structural obstruction
(Remark~\ref{rem:structural-obstruction}) shows that algebraic
constraints from the MC equation act on the real spectral line,
while the zeta zeros live at complex spectral parameters.  No
amount of clever identification can bridge this gap: it is the
difference between a finite-dimensional algebraic structure
(the shadow tower, determined by three numbers $\kappa$, $\alpha$,
$S_4$) and an infinite-dimensional analytic object (the zeta
function, with infinitely many independent zeros).

The shadow tower \emph{detects} arithmetic---Hecke eigenforms
for lattice VOAs, quadratic fields for all VOAs at rational $c$---but
it does not \emph{control} the arithmetic.  The detection is a
theorem (Theorem~\ref{thm:shadow-spectral-correspondence}).  The
control would require a mechanism for the MC equation to force
zeros onto critical lines, and no such mechanism exists.

What remains is beautiful and proved.  The bar construction of a
chiral algebra is a Fourier transform that simultaneously produces
topological invariants (Mumford classes), arithmetic invariants
(quadratic fields), and algebraic invariants (Koszul duals) from
a single MC element.  The gravitational shadow tower, specialised
to rational central charge, assigns a number field to each minimal
model.  The Ising model lives over $\Q(\sqrt{5})$.  The self-dual
gravity at $c = 13$ lives over $\Q(\sqrt{870})$.  Whether these
number fields have structure beyond what is visible in the
specialisation table is an open question at the intersection of
vertex algebra theory and algebraic number theory.

The bar construction makes the question precise.  It does not
answer it.


% =================================================================
\section{Session notes: 1 April 2026}
\label{sec:session-2026-04-01}
% =================================================================

Six findings from the session.  The first is a theorem; the
remaining five are its consequences.


% -----------------------------------------------------------------
\subsection{The DS-HPL transfer theorem}
\label{sec:ds-hpl-transfer}
% -----------------------------------------------------------------

The single strongest new result.

\begin{theorem}[DS-HPL transfer; proved]
\label{thm:ds-hpl-transfer-wn}
Let $\fg$ be a simple Lie algebra, $k$ a non-critical level, and
$\cW_k(\fg)$ the principal Drinfeld--Sokolov reduction of
$V_k(\fg)$.  The homological perturbation lemma transfers the
dg-shifted Yangian datum $(m, r, \Delta)$ from the affine side to
the $\cW$-algebra side via the BRST strong deformation retract
$(\iota, p, h)$ with $\mathrm{gh}(h) = -1$.  The transferred
coproduct $\Delta^W$ is strictly primitive at every arity
(equation~\eqref{eq:gravity-primitive-wn}), with no corrections.
\end{theorem}

\begin{proof}
The ghost-number argument of
Theorem~\ref{thm:gravitational-primitivity-wn}
(\S\ref{sec:primitive-coproduct}) applies verbatim: every HPL
tree with an internal vertex shifts ghost number away from~$0$,
and $p \otimes p$ annihilates the result.
\end{proof}

The theorem says that Drinfeld--Sokolov reduction kills the
coproduct channel completely.  Only the $r$-matrix channel
survives.  The $\cW$-algebra inherits the $\Ainf$ products
$m_k^W$ (infinitely many, for class~M) and the collision
residue $r^W(z)$, but its coproduct is as trivial as a
one-dimensional algebra's.  It produces nothing.  It only
braids.

\sep

A remark on what is transferred and what is not.  The SDR
$(\iota, p, h)$ transfers three structures simultaneously:
\begin{center}
\small
\renewcommand{\arraystretch}{1.15}
\begin{tabular}{@{}lll@{}}
\textbf{Structure} & \textbf{Transferred result} & \textbf{Status} \\
\hline
$\Ainf$ products $m_k$ & nontrivial for all $k \geq 2$
  & class~M \\[3pt]
$r$-matrix $r(z)$ & nontrivial, all poles survive
  & quasi-trigonometric \\[3pt]
Coproduct $\Delta_z$ & primitive (all corrections vanish)
  & proved
\end{tabular}
\end{center}

\medskip

The $\Ainf$ products are nontrivial because the perturbation
$\delta_{\mathrm{BRST}}$ has ghost-number-preserving components
that survive the projection $p$.  The $r$-matrix is nontrivial
because it arises from the collision residue, which is a
genus-$0$ operation with no homotopy insertions at the output.
The coproduct is primitive because it is a genus-$0$,
arity-$2$, \textit{splitting} operation: the output has two
tensor factors, and every homotopy insertion shifts ghost number
into the range that $p \otimes p$ annihilates.

The asymmetry is that collision (computing $r(z)$) applies $p$
once, while splitting (computing $\Delta_z$) applies
$p \otimes p$: the tensor product doubles the projection's
power to kill ghost-shifted outputs.


% -----------------------------------------------------------------
\subsection{Gravitational primitivity as physical principle}
\label{sec:gravitational-primitivity-physical}
% -----------------------------------------------------------------

The physical content of gravitational primitivity
(\S\ref{sec:primitive-coproduct}): in three dimensions, there are
no local gravitational degrees of freedom and no propagating
gravitons.  The coproduct channel is trivial
(equation~\eqref{eq:gravity-primitive-wn}), and only the
$r$-matrix channel survives.  BTZ black holes are monodromies of
the $r$-matrix braiding, not vertices of the coproduct.

\sep

The $(m, \Delta)$-complexity plane organises the five standard
HT theories:

\begin{center}
\small
\renewcommand{\arraystretch}{1.15}
\begin{tabular}{@{}lcc@{}}
\textbf{Theory} & $m$-depth & $\Delta$-depth \\
\hline
Chern--Simons & $2$ (quadratic) & $\geq 1$ (nontrivial) \\[3pt]
Twisted holography & $2$ & $\geq 1$ \\[3pt]
M2 & $2$ & $\geq 1$ \\[3pt]
Gravity & $\infty$ (class M) & $0$ (primitive) \\[3pt]
M5 & $\infty$ & $0$
\end{tabular}
\end{center}

\medskip

\noindent
Gravity sits at $(\infty, 0)$: maximal $\Ainf$ complexity,
vanishing coproduct complexity.  Gauge theory sits at
$(2, \geq 1)$: minimal $\Ainf$ complexity, nontrivial
coproduct.  DS reduction moves a theory from the gauge corner
to the gravity corner by killing the coproduct and inflating
the $\Ainf$ tower.


% -----------------------------------------------------------------
\subsection{Five worked examples as $\Theta^{\mathrm{oc}}$
  instantiations}
\label{sec:five-oc-instantiations}
% -----------------------------------------------------------------

The completed eight-fold platonic datum
$\Pi^{\mathrm{oc}}_X(\cA)$ for the five canonical HT systems:

\medskip

\noindent
\textbf{1.\ Chern--Simons.}\enspace
$\cA = V_k(\fg)$.  Class~L, shadow depth $3$.
$\kappa + \kappa^! = 0$ (Feigin--Frenkel).  Coproduct:
nontrivial (Casimir cross-term).  Sewing: determinant.  The
entire $\Ainf$ structure is formal: $m_k^H = 0$ for $k \geq 3$ on cohomology.
The platonic package is finite-dimensional at each arity.

\medskip

\noindent
\textbf{2.\ Three-dimensional gravity.}\enspace
$\cA = \mathrm{Vir}_c$.  Class~M, shadow depth $\infty$.
$\kappa + \kappa^! = 13 \neq 0$.  Coproduct: \textbf{primitive}
(ghost-number obstruction from DS).  Sewing:
Pfaffian-squared.  The quartic contact invariant
$Q^{\mathrm{ct}} = 10/[c(5c+22)]$ controls the entire infinite
tower.  The $r$-matrix
$r(z) = (c/2)/z^3 + 2T/z$ carries all scattering.

\medskip

\noindent
\textbf{3.\ Twisted holography.}\enspace
$\cA = V_k(\fg)$ (affine Kac--Moody, Costello--Gaiotto).
Class~L, shadow depth $3$.  Multi-channel: for $\cW_3$, the
genus-$2$ cross-channel contribution is
\[
F_2^{\mathrm{cross}} \;=\; \frac{c + 120}{16c}.
\]
The holographic dictionary is the Koszul triangle read through
Theorems~A--H.

\medskip

\noindent
\textbf{4.\ M2-brane.}\enspace
$\cA = \mathrm{DDCA}_k(\mathfrak{gl}_K)$.  Class~L.  Coproduct:
nontrivial.  Two-variable $r$-matrix $r(z_1, z_2)$.  No DS
reduction: the $\cW$-algebra structure is absent.  This is
gauge theory, not gravity, and the nontrivial coproduct
confirms it.

\medskip

\noindent
\textbf{5.\ M5-brane / $6$d $(2,0)$.}\enspace
$\cA = \cW_\infty(K)$ at large~$N$.  Class~M, shadow depth
$\infty$.  Coproduct: \textbf{primitive} (same ghost-number
argument as Virasoro, since $\cW_N$ arises by DS reduction).
Per-spin $r$-matrix.  Shadow depth increases under DS.

\sep

The shadow dichotomy across the five theories:

\begin{center}
\small
\renewcommand{\arraystretch}{1.15}
\begin{tabular}{@{}lcl@{}}
\textbf{Lineage} & \textbf{Class} & \textbf{Mechanism} \\
\hline
Current-algebra (CS, M2) & L & finite towers, no DS \\[3pt]
$\cW$-algebra (gravity, $\cW_3$, M5) & M & infinite towers,
  DS origin
\end{tabular}
\end{center}

\medskip

DS is the shadow-depth escalator.  It converts class~L (finite,
formal) to class~M (infinite, genuinely $\Ainf$).  The
ghost-number filtration that kills the coproduct is the same
filtration that inflates the $\Ainf$ tower.  The two effects
are a single mechanism.


% -----------------------------------------------------------------
\subsection{The open-sector primitive}
\label{sec:open-sector-primitive-session}
% -----------------------------------------------------------------

The primitive object is $\cC_{\mathrm{op}}$, not $\cA$ or
$\barB(\cA)$.

For Chern--Simons: $\cC_{\mathrm{op}} = V_k(\fg)\text{-mod}$,
the category of $V_k(\fg)$-modules as a factorisation
dg-category with meromorphic tensor product.  The boundary
algebra $A_b = \End(\mathrm{trivial})$ recovers $V_k(\fg)$.
The chiral derived center $\Zder = C^\bullet_{\mathrm{ch}}(\cA_b,
\cA_b)$ classifies bulk operators.

Modularity is trace plus clutching on $\cC_{\mathrm{op}}$, not
an axiom on the closed algebra.  The annulus trace
\[
\mathrm{Tr}\colon K_0(\cC_{\mathrm{op}})
\;\longrightarrow\;
Z_1(\cA)
\]
gives the genus-$1$ shadow.  The full genus-$g$ amplitude is the
composition of $g$ clutching maps on $\cC_{\mathrm{op}}$.  The
modular $S$-matrix is a derived property of the trace, not an
input to the theory.

The boundary algebra $A_b$ is a chart: a Morita-dependent
presentation of $\cC_{\mathrm{op}}$.  Different compact
generators produce Morita-equivalent charts.  The category is
the invariant; the algebra is a coordinate system.


% -----------------------------------------------------------------
\subsection{BV/BRST $=$ bar: honest status}
\label{sec:bv-brst-bar-honest}
% -----------------------------------------------------------------

Three levels of identification.  The first is proved; the
second is conditionally proved; the third is conjectural.

\begin{center}
\small
\renewcommand{\arraystretch}{1.15}
\begin{tabular}{@{}llc@{}}
\textbf{Genus} & \textbf{Claim} & \textbf{Status} \\
\hline
$g = 0$ & algebraic BRST $=$ bar differential & proved \\[3pt]
$g = 0$, higher arity
  & Feynman-diagrammatic $m_k$ $=$ bar $m_k$
  & proved \\[3pt]
$g \geq 1$ & BV quantum master equation $=$ curved bar
  & conjectural
\end{tabular}
\end{center}

\medskip

The bar complex IS well-defined at all genera: the inductive
genus determination from Volume~I combined with the HS-sewing
theorem produces bar amplitudes at every genus.  What is open
is whether these bar amplitudes coincide with the BV-BRST
partition function at genus $g \geq 1$.  This is
\texttt{conj:master-bv-brst}.

The identification at genus~$0$ is an algebraic comparison:
the BRST differential $Q = \partial + \partial^{\mathrm{BRST}}$,
restricted to the physical subspace, reproduces the bar
differential $d_{\barB}$.  The identification at genus~$1$
would require matching the BV quantum master equation
$\hbar\Delta_{\mathrm{BV}}S + \tfrac{1}{2}\{S, S\} = 0$ with
the curved bar relation $\dfib^{\,2} = \kappa \cdot \omega_1$.
No obstruction to this matching is known; no proof exists.

This session corrected four locations in the manuscript that
overclaimed genus-$1$ BV/BRST $=$ bar as proved.  The bar
complex is the proved object.  The BV-BRST comparison is the
conjectural bridge.


% -----------------------------------------------------------------
\subsection{The standalone paper}
\label{sec:standalone-paper}
% -----------------------------------------------------------------

\textit{Shadow Towers and Algebraicity}, $25$ pages, eight
sections.

\begin{enumerate}[label=\S\arabic*.,nosep]
\item Introduction and motivation.  The bar construction as
  categorical logarithm; the shadow tower as its Taylor
  expansion.
\item The MC algebra.  The modular cyclic deformation complex
  $\mathrm{Def}^{\mathrm{mod}}_{\mathrm{cyc}}(\cA)$ and the
  universal element $\Theta_\cA$.
\item The shadow tower.  Projection to the primary line;
  recursive computation of $S_r$ from the MC equation.
\item The Riccati algebraicity theorem.  The MC equation on a
  single primary line reduces to a Riccati ODE whose solution
  is $H(t) = t^2\sqrt{Q_L(t)}$, algebraic of degree~$2$.
  Full proof.
\item Four-class classification.  The critical discriminant
  $\Delta = 8\kappa S_4$ partitions modular Koszul algebras
  into four classes:
  \begin{enumerate}[label=(\alph*),nosep]
  \item \textbf{G} (Gaussian): $r_{\max} = 2$.  Heisenberg.
  \item \textbf{L} (Lie/tree): $r_{\max} = 3$.  Affine
    Kac--Moody.
  \item \textbf{C} (contact/quartic): $r_{\max} = 4$.
    $\beta\gamma$.
  \item \textbf{M} (mixed): $r_{\max} = \infty$.  Virasoro,
    $\cW_N$.
  \end{enumerate}
\item Explicit shadow towers.  Complete computation for the
  standard landscape: Heisenberg, affine $\fsl_2$,
  $\beta\gamma$, Virasoro, $\cW_3$.  Tables of $S_r$ through
  arity~$8$.
\item $\cW_3$ cross-channel genus-$2$ computation.
  Multi-generator shadow tower with two primary lines
  ($T$ and $W$), propagator variance $\delta_{\mathrm{mix}}$,
  genus-$2$ cross-channel formula.
\item Outlook.  Pixton conjecture (class~M algebras generate
  the Pixton ideal via infinite MC tower), analytic completion
  programme, non-standard families.
\end{enumerate}

The paper extracts from Volume~I the self-contained story of
the shadow tower, from its definition through the algebraicity
theorem to the four-class classification and the explicit
standard-landscape computations.  It is the first accessible
entry point to the monograph's nonlinear characteristic layer.


% =================================================================
\section{The graviton is composite}
\label{sec:graviton-composite}
% =================================================================

The Sugawara stress tensor
\[
T \;=\; \frac{1}{2(k+h^\vee)}\sum_a \normord{J^a J_a}
\]
is quadratic in affine currents.  It is not a fundamental field.
The graviton is composite.

The BRST complex of the Drinfeld--Sokolov reduction remembers
this: the transferred coproduct is strictly primitive
(Theorem~\ref{thm:gravitational-primitivity-wn},
equation~\eqref{eq:gravity-primitive-wn}).  No cross-term, no
splitting, at every arity and every principal DS reduction.

\sep

In gauge theory: the gluon $J^a$ IS the connection.
Fundamental.  Ghost number~$0$.  No passage through~$h$.  The
BRST homotopy does not obstruct the transferred coproduct.
The Casimir cross-term survives projection:
\[
\Delta_z(T_{\mathrm{aff}})
\;=\;
\tau_z(T_{\mathrm{aff}}) \otimes 1
\;+\; 1 \otimes T_{\mathrm{aff}}
\;+\; \frac{1}{k+h^\vee}\,\Omega(w{+}z,\, w),
\]
where $\Omega = \sum_a J^a \otimes J_a$ is the split Casimir.
A production vertex.  A gluon bound state dissociates.

In gravity: the graviton $T$ IS the curvature.  Composite.
Ghost-burdened.  Every HPL tree at ghost number~$0$ dies under
projection.  Gravitons do not split.  They braid.

\sep

The $r$-matrix $r^{\mathrm{Vir}}(z) = (c/2)/z^3 + 2T/z$
carries all gravitational scattering.  The cubic pole: the bar
kernel $d\log(z-w)$ absorbs one power from the quartic OPE.
The Virasoro OPE has poles at $z^{-4}$, $z^{-2}$, $z^{-1}$.
The $r$-matrix has poles at $z^{-3}$, $z^{-1}$.  The
logarithmic derivative is always one degree milder than the
function itself.  This absorption is not an approximation.  It
is the structural content of calling the bar construction a
categorical logarithm.

\sep

The five-theory comparison (see \S\ref{sec:five-theories})
confirms the pattern: DS reduction creates depth (L $\to$ M)
and kills the coproduct.  The distinction is structural
(the ghost-number filtration), not parametric (the central charge).

\sep

The bar complex of a chiral algebra knows more about quantum
gravity than any perturbative expansion.

It knows the exact genus-$g$ free energy at every genus, in
closed form:
$F_g = (c/2) \cdot \lambda_g^{\mathrm{FP}}$, the
$\hat{A}$-genus coefficients times the modular characteristic.

It knows the entire infinite tower of higher-arity corrections,
as the Taylor expansion of an algebraic function of degree~$2$:
$H(t) = t^2\sqrt{Q_{\mathrm{Vir}}(t)}$, three numbers
$(\kappa, \alpha, S_4)$, one quadratic form, one square root.

It knows that the gravitational coproduct is primitive
(\S\ref{sec:primitive-coproduct}): black holes do not split.

It knows the arithmetic.  At rational central charge $c_0$,
the shadow metric $Q_{\mathrm{Vir}}$ specialises to a specific
quadratic polynomial over~$\Q$, and
$K(c_0) = \Q(\sqrt{d(c_0)})$ is a specific imaginary quadratic
number field.  The Bekenstein--Hawking entropy $S = 2\pi\sqrt{ch/6}$
is the leading coefficient.  The discriminant
$-320c^2/(5c+22)$ classifies the depth.  The class number of
$K(c_0)$ measures the arithmetic complexity of gravity at that
central charge.  The Ising model lives over $\Q(\sqrt{-10})$.
The self-dual gravity at $c = 13$ lives over $\Q(\sqrt{-870})$,
and the complementarity constant $6960 = 2^4 \cdot 3 \cdot 5
\cdot 29$ appears as $-80 \cdot 87 = -6960$.

\sep

The bar complex of a chiral algebra is a $1$-form, a square
root, and a ghost.

The $1$-form $\eta = d\log(z_i - z_j)$ extracts the OPE.
Arnold's relation makes it nilpotent.  The genus-completed
differential $D_\cA$ satisfies $D_\cA^2 = 0$, and
$\Theta_\cA = D_\cA - d_0$ is the universal MC element.

The square root $\sqrt{Q_L}$ generates the tower.  One
quadratic form, one algebraic function, infinitely many
rational invariants.  The depth classification, the growth
rate, the convergence behaviour, the Galois theory, the Epstein
zeta function, the discriminant complementarity: all
are properties of this single square root.

The ghost kills the splitting
(Theorem~\ref{thm:gravitational-primitivity-wn}): the BRST
homotopy carries ghost number $\pm 1$, the projection
annihilates it, and the coproduct vanishes.

Everything else is a projection of these three.


% ============================================================
\section{Koszul--Epstein and quantum gravity}
\label{sec:koszul-epstein-gravity}
% ============================================================

\subsection{The arithmetic content of the gravitational bar complex}

The bar complex of the Virasoro algebra at central charge~$c$
carries three layers of arithmetic data.

\textbf{Layer 1: the Heisenberg core.}
The Miura map $\hat\Upsilon\colon W_\kappa(\mathfrak{sl}_2) \to
V^{\mathfrak{h}_\kappa}(\mathfrak{h})$ embeds the Virasoro
algebra into a rank-$1$ Heisenberg.  At the level of the
sewing determinant, this produces an exact Euler product:
\[
\sum_{N \ge 1} \sigma_{-1}(N)\,q^N
\;=\; -\log\bigl(\eta(q)^2\bigr) - \tfrac{\tau}{12},
\qquad
\sum_{N \ge 1} \sigma_{-1}(N)\,N^{-s}
\;=\; \zeta(s)\,\zeta(s+1).
\]
The Heisenberg core is arithmetically transparent: all
nontrivial genus-one arithmetic is a consequence of the
divisor function~$\sigma_{-1}$ being multiplicative.

\textbf{Layer 2: the Miura defect.}
The vacuum character of~$W_N$ splits as
$\chi^{\mathrm{vac}}_{W_N}(q)
= \chi_H(q)^{N-1} \cdot
\prod_{m=1}^{N-1}(1-q^m)^{N-m}$.
All nontrivial genus-one arithmetic beyond the Heisenberg core
sits in the finite Miura defect factor
$D_N(q) = \prod_{m=1}^{N-1}(1-q^m)^{N-m}$.
For~$N = 2$ (Virasoro): $D_2(q) = 1 - q$, removing the weight-$1$
mode.

\textbf{Layer 3: the shadow tower.}
The infinite shadow tower $\{S_r(c)\}_{r \ge 2}$, determined by
$(\kappa, \alpha, S_4)$ via the MC recursion, encodes the full
$A_\infty$ deformation data.  Its spectral measure~$\rho$ is
uniquely determined by the Carleman condition.  The MC recursion
propagates Hecke equivariance from the character level to all
arities (Route~C, proved for the standard landscape).

\subsection{The gravitational dichotomy and the Koszul--Epstein frontier}

The gravity dichotomy at $c = 26$ concerns the vanishing of
$\kappa_{\mathrm{eff}} = \kappa(\mathrm{matter})
+ \kappa(\mathrm{ghost})$.  The Koszul--Epstein question
concerns the critical-line properties of the arithmetic
shadow.  These are unrelated:

\begin{itemize}
\item $\kappa_{\mathrm{eff}} = 0$ at $c = 26$ determines whether
  the genus tower terminates (the Clifford-to-exterior
  degeneration of the transgression algebra).

\item The critical-line question concerns the zeros of
  $\zeta(s)$ appearing in the Heisenberg sewing determinant.
  These zeros are independent of the central charge: the sewing
  lift $S_{\mathrm{Vir}}(u) = \zeta(u+1)(\zeta(u) - 1)$ is the
  same function for all~$c$.
\end{itemize}

The Koszul self-duality point $c = 13$, where
$\kappa + \kappa' = 13$ and $\delta_\kappa = 0$, is the point
where Verdier--Hecke commutation forces the constrained Epstein
zeta to factor into standard $L$-functions.  This is a genuine
arithmetic structure of the self-dual gravitational theory.
It does not, however, produce a critical-line theorem.

\subsection{Honest assessment}

The MC element~$\Theta_\cA$ of the gravitational bar complex
produces genuine arithmetic structure: the Carleman-determined
spectral measure, Route~C prime-locality, the shadow--spectral
correspondence for lattice VOAs, the Verdier--Hecke factorisation
at self-dual points.  These are proved theorems.

The programme does not produce a critical-line theorem.  The
structural obstruction is that the shadow spectral measure
and the automorphic spectral measure are different objects, and
the map between them (the Rankin--Selberg integral) is not
sign-preserving.  No mechanism for positivity has been
identified.  The bar--cobar homotopy equivalence provides
deformation-theoretic transport, not arithmetic positivity.

The surviving content: the quartic resonance class
$R_4^{\mathrm{mod}}(\cA)$ is the first nonlinear place where
Koszul duality constrains the automorphic spectral decomposition
beyond the character level.  Whether this constraint carries
critical-line information is open.  The correct next computation
is the quartic Gram-side matching with the crossing-weighted
bilinear residue kernel (using actual DOZZ/Toda structure
constants), not the bare spectral measure (which fails by four
to five orders of magnitude).

The bar complex of a chiral algebra knows the exact genus tower,
the arithmetic of the shadow, and the depth classification.
It does not know the zeros of~$\zeta(s)$.


% ============================================================
\section{Spectrum and gravity}
\label{sec:spectrum-gravity}
% ============================================================

The gravitational bar complex carries three kinds of data:
the algebraic engine (Theorems A--H, the MC
element~$\Theta_\cA$), the arithmetic shadow (the shadow tower,
the depth classification, the Epstein zeta), and the spectral
content (the primary spectrum, the partition function, the
constrained Epstein series~$\varepsilon^c_s$).  These three
overlap but do not coincide.  This section records the definitive
assessment of their relationship in the gravitational context.

\subsection{The spectrum as a categorical invariant}

The scalar primary spectrum
$\mathrm{Spec}_{\mathrm{sc}}(\cA) = \{(\Delta_i, m_i)\}$ of a
2d~CFT at central charge~$c$ is an invariant of the
\emph{modular tensor category}
$\cA\text{-}\mathrm{mod}^{\mathrm{ss}}$ together with the
modular-invariant partition function $Z(\tau,\bar\tau)
= \sum_{i,j} M_{ij}\,\chi_i(\tau)\,\overline{\chi_j(\tau)}$.
The multiplicity matrix~$M_{ij}$ encodes the physical spectrum;
the characters $\chi_i$ encode the algebraic spectrum.

The bar complex $B(\cA)$ is built from the algebra~$\cA$
(specifically, from the vacuum module and its OPE), not from
the category of modules.  The bar-cobar machine determines the
Koszul dual~$\cA^!$, the MC element~$\Theta_\cA$, and the shadow
tower~$\{S_r\}$.  It does \emph{not} determine the modular
$S$-matrix or the multiplicity matrix.  The primary spectrum is
therefore external to the bar construction.

For the gravitational theory (Virasoro at central charge~$c$),
the vacuum module is the unique irreducible highest-weight module
$L(c,0)$, and the partition function $Z_{\mathrm{grav}}
= |\chi_0(\tau)|^2 + \sum_{h > 0} m_h\,|\chi_h(\tau)|^2$ depends
on the dynamical data $\{m_h\}_{h > 0}$: the spectrum of
primary states.  The bar complex of~$\mathrm{Vir}_c$ sees the
algebra, not the specific gravity theory built from it.

\subsection{Which constructions partially capture the spectrum}

Five constructions each capture a shadow of the gravitational
primary spectrum.  None captures it fully from within the bar
machine.

\begin{enumerate}
\item \emph{Zhu's algebra $A(\mathrm{Vir}_c)$.}
  For the universal Virasoro: $A(V_c) \cong \C[x]$.  The simple
  modules of $\C[x]$ correspond to points $x = h \in \C$: the
  conformal dimensions.  Zhu's algebra determines the
  \emph{set} of allowed lowest weights $\{h_{r,s}\}$ for the
  minimal model at $c = c_{p,q}$ (via the quotient $A(L(c,0))
  = \C[x]/(f_c(x))$).  But it does not determine the partition
  function multiplicities.  At generic~$c$ (irrational, or
  rational non-minimal), $A(V_c) = \C[x]$ admits a continuum
  of simple modules, and no finite selection rule from $A(V_c)$
  picks out the physical spectrum.

\item \emph{Associated variety $X_{\mathrm{Vir}_c}$.}
  For Virasoro: $X_{V_c} = \C$ (one-dimensional, since Virasoro
  has a single strong generator).  The dimension of $X_V$
  controls the asymptotic growth of the character and hence the
  leading pole of the constrained Epstein series at $s = c/2$.
  This is the Cardy regime: $\log \rho(E) \sim
  2\pi\sqrt{c\,E/3}$ for large conformal dimension~$E$.

\item \emph{Bar cohomology $H^*(B(\mathrm{Vir}_c))$.}
  Concentrated in bar degree~$1$ (Koszulness).  The graded
  dimensions $\dim H^1(B(\mathrm{Vir}))_h$ count generators of
  the Koszul dual $\mathrm{Vir}^!_{26-c}$ at weight~$h$.  These
  are the \emph{Motzkin differences} $M(h+1) - M(h)$, supported
  only at even weights, independent of~$c$, and growing
  as~$3^h$.  The quasi-primary count (partitions minus
  descendants) is a different sequence, nonzero at odd weights
  and growing as $e^{\pi\sqrt{2h/3}}$.  Bar cohomology sees the
  dual algebra, not the spectrum.

\item \emph{Shadow tower $\{S_r\}$.}
  The shadow tower extracts three numbers $(\kappa, \alpha, S_4)$
  from the OPE and packages them into the shadow metric~$Q_L$.
  The shadow-metric Epstein zeta $\varepsilon_{Q_L}(s)$ is a
  binary quadratic form sum, unrelated to the primary spectrum
  except through a chain of projections.  For lattice VOAs, the
  shadow--spectral correspondence recovers the full $L$-function
  package; for Virasoro, the correspondence is incomplete (the
  shadow has $3$~parameters; the spectrum has infinitely many).

\item \emph{The chiral derived centre
  $\cZ^{\mathrm{der}}_{\mathrm{ch}}(\cA)$.}
  The universal bulk algebra, built as chiral Hochschild cochains
  $C^\bullet_{\mathrm{ch}}(\cA_b, \cA_b)$, classifies bulk
  operators.  Its cohomology controls the BRST-closed observables
  of the 3d HT theory.  But it is a genus-$0$ invariant of the
  open/closed Swiss-cheese structure, and does not directly encode
  the genus-$1$ partition function.  The trace on the derived
  centre (the annulus trace~$\mathrm{Tr}_\cA$) recovers the
  Hochschild homology $\mathrm{HH}_*(\cA)$, which is the first
  modular shadow of the open sector.  This is related to, but
  not identical with, the partition function.
\end{enumerate}

\subsection{The gravitational spectral gap}

In the gravitational context, the primary spectrum has additional
structure: the conformal bootstrap imposes crossing-symmetry
constraints on the four-point function, which restrict the
allowed $\{(\Delta_i, m_i)\}$ beyond what the chiral algebra
alone determines.  The Benjamin--Chang constrained Epstein
series $\varepsilon^c_s$ is the Eisenstein spectral coefficient
of $\hat Z^c$, and its functional equation involves
$\zeta(2s)/\zeta(2s-1)$, introducing poles at
$s = (1+\rho_n)/2$ where $\rho_n$ are nontrivial zeros
of~$\zeta$.

The pole \emph{locations} are universal: they come from
$\zeta(2s-1) = 0$ in the $\Gamma$-factor ratio $F_c(s)$, and
do not depend on the specific CFT.  The pole \emph{residues}
are CFT-specific: they are integrals of the primary-counting
function against residues of Eisenstein series on the modular
fundamental domain.  The bar complex constrains neither the
locations (which are number-theoretic: the zeros of~$\zeta$)
nor the residues (which are spectral: the primary multiplicities
$m_h$).

For the gravitational theory at $c = 26$ (critical string):
$\kappa_{\mathrm{eff}} = \kappa(\mathrm{matter})
+ \kappa(\mathrm{ghost}) = 0$, and the genus tower terminates.
But the primary spectrum is still infinite (the Virasoro module
at $c = 26$ has infinitely many primaries at all weights
$h \geq 2$), and $\varepsilon^{26}_s$ is a nontrivial Dirichlet
series.  The genus tower's termination is a statement about the
$\hat A$-genus factor~$\kappa$; it says nothing about the
spectral content.

For the gravitational theory at $c = 13$ (self-dual point):
Verdier--Hecke commutation forces the constrained Epstein to
factor into standard $L$-functions.  This is the strongest
structural result: self-duality of the chiral algebra produces
a factorisation of the spectral invariant.  But even at
$c = 13$, the factorisation does not determine the zeros of the
constituent $L$-functions.

\subsection{Is $\varepsilon^c_s$ a bar-complex invariant?}

\begin{proposition}[The constrained Epstein is not a bar-complex invariant]
\label{prop:constrained-epstein-not-bar}
The constrained Epstein series $\varepsilon^c_s(\cA)$ cannot be
reconstructed from the bar complex $B(\cA)$ and its cohomology.
\end{proposition}

\begin{proof}
Three independent arguments.

\textbf{(a) Categorical argument.}
Two non-isomorphic modular-invariant partition functions for the
same chiral algebra produce different constrained Epstein series.
Example: the Ising model ($c = 1/2$) has the diagonal partition
function $Z_{\mathrm{diag}} = |\chi_0|^2 + |\chi_{1/2}|^2
+ |\chi_{1/16}|^2$ and the charge-conjugation partition function
(which is the same, since Ising is self-conjugate).  But for
$\widehat{\mathfrak{su}(2)}_k$ at $k \geq 10$, there exist
exceptional modular invariants (the $D$-type and $E$-type
classifications of Cappelli--Itzykson--Zuber) that produce
different primary spectra from the same chiral algebra.  The
bar complex $B(\hat\fg_k)$ is the same for all choices of
modular invariant; the constrained Epstein series differs.

\textbf{(b) Dimensional argument.}
The bar complex is determined by the OPE data of the vacuum
module (a countable set of structure constants).  The primary
spectrum is determined by the modular invariant \emph{and} the
full representation theory (the braided tensor structure, the
modular $S$-matrix).  The latter is strictly richer: two algebras
with the same OPE but different categories of modules (e.g.,
$V_k(\fg)$ at the same $k$ but with different conformal
extensions) have the same bar complex but different spectra.

\textbf{(c) Weight argument.}
$\dim H^1(B(\mathrm{Vir}))_h$ at odd $h$ vanishes (the Koszul
dual has generators only at even weights), while the quasi-primary
count at odd $h$ is generically nonzero.  No Dirichlet series
built from bar cohomology dimensions can reproduce
$\varepsilon^c_s$, which sums over all primary weights including
odd ones.
\end{proof}

\subsection{The honest three-layer picture}

The gravitational bar complex produces three layers of arithmetic
data, each genuine and proved.

\textbf{Layer 1: the genus expansion.}
$F_g(\cA) = \kappa(\cA) \cdot \lambda_g^{\mathrm{FP}}$
for uniform-weight algebras.  This is the leading scalar projection
of the MC element $\Theta_\cA$, generated by $\hat A(i\hbar) - 1$.
It captures \emph{one number} ($\kappa$) from the full partition
function.

\textbf{Layer 2: the shadow tower.}
$\{S_r\}_{r \geq 2}$, determined by $(\kappa, \alpha, S_4)$.
This captures \emph{three numbers} from the OPE, packaged into
the shadow metric~$Q_L$ and the quadratic extension~$K_L$.  The
shadow--spectral correspondence for lattice VOAs, Route~C
prime-locality, and the Verdier--Hecke factorisation are genuine
arithmetic results at this layer.

\textbf{Layer 3: the spectral content.}
$\varepsilon^c_s(\cA) = \sum (2\Delta)^{-s}$, encoding the full
scalar primary spectrum.  This layer is external to the bar
complex: it requires the partition function, which is a
genus-$1$ categorical invariant, not a genus-$0$ algebraic one.

The gap between Layers 2 and 3 is the gap between the
shadow tower (genus-$0$, $3$~parameters) and the primary spectrum
(genus-$1$, $\infty$~parameters).  No mechanism in the MC equation
or the bar-cobar machine bridges this gap: the shadow tower is
determined by finitely many OPE coefficients, while the spectrum
encodes infinitely many independent quantities.

The honest statement for the gravitational theory:
\emph{the bar complex of the Virasoro algebra at central
charge~$c$ determines the exact genus expansion
$F_g = (c/2)\,\lambda_g^{\mathrm{FP}}$, the infinite shadow
tower with growth rate~$\rho(c)$ and discriminant
$\Delta = 80/(5c{+}22)$, and the four-class depth assignment
(class~M for all $c \neq 0$).  It does not determine the
primary spectrum, the partition function, or the constrained
Epstein series.  The spectrum is a representation-theoretic
invariant; the bar complex is an algebra-theoretic invariant.
They live in different categories.}


% ============================================================
\section{Derived global sections in quantum gravity}
\label{sec:derived-global-sections-qg}
% ============================================================

The gravitational bar complex defines a sheaf on moduli space
whose derived global sections are algebraic invariants.  This
section records the definitive assessment of what these
algebraic invariants capture in the gravitational context,
correcting the false claim that ``the bar complex cannot see the
spectrum because integration over moduli space is analytic.''

\subsection{The sheaf-theoretic framework}

For Virasoro at central charge~$c$, the bar complex at
genus~$g$ defines a sheaf of dg modules on
$\overline{\cM}_g$.  The key objects:

\begin{enumerate}
\item The \emph{virtual bar $K$-class}
  $V_{\mathrm{Vir}_c} := [R\pi_{g*}\,\barB^{(g)}(\mathrm{Vir}_c)]
  \in K^0(\overline{\cM}_g)$, computed by derived pushforward
  along the universal curve.  This is well-defined: the
  perfectness criterion (Koszul acyclicity + finite-dimensional
  fiber cohomology) is satisfied for all modular Koszul algebras.

\item The \emph{center sheaf}
  $\cZ_{\mathrm{Vir}_c} := R^*\pi_{g*}(\barB^{(g)},\, d_{\mathrm{fib}})$,
  a local system on $\overline{\cM}_g$ with fiber
  $H^{\mathrm{ch}}_*(\Sigma_g, \mathrm{Vir}_c)$.

\item The \emph{ambient complex}
  $C_g(\mathrm{Vir}_c) = R\Gamma(\overline{\cM}_g,\,
  \cZ_{\mathrm{Vir}_c})$, the derived global sections.
  This is a purely algebraic operation on a constructible sheaf
  over a proper DM stack.

\item The \emph{GRR hierarchy}:
  \begin{align*}
  c_1(\det V_{\mathrm{Vir}_c}) &= \tfrac{c}{2} \cdot \lambda, \\
  \mathrm{ch}(V_{\mathrm{Vir}_c}) &\leftrightarrow \Delta_{\mathrm{Vir}}(x)
    \quad\text{(Motzkin differences)}, \\
  \mathrm{Hol}(V_{\mathrm{Vir}_c}) &= \Pi_{\mathrm{Vir}_c}
    \quad\text{(MCG monodromy)}.
  \end{align*}
  All three levels are algebraic.
\end{enumerate}

So the bar complex \emph{does} produce algebraic invariants
from integration over moduli space.  The claim that this
integration is ``analytic, not algebraic'' was a false idea.
Derived global sections, Chern character, GRR pushforward,
and monodromy representation are all algebraic operations.

\subsection{The gravitational spectral gap: precise diagnosis}

The GRR hierarchy captures the \emph{algebra's} modular
structure.  It does not capture the \emph{gravitational
spectrum}.  The precise diagnosis:

\textbf{(i) The GRR scalar level} recovers $\kappa(\mathrm{Vir}_c)
= c/2$.  This is one number.  The gravitational partition function
$Z(\tau,\bar\tau) = |\chi_0(\tau)|^2 + \sum_{h > 0} m_h\,
|\chi_h(\tau)|^2$ encodes infinitely many numbers $\{m_h\}$.

\textbf{(ii) The GRR spectral level} recovers the bar cohomology
dimensions $\dim H^1(B(\mathrm{Vir}))_h$, which are the Motzkin
differences $M(h+1) - M(h)$.  These count generators of the
Koszul dual $\mathrm{Vir}^!_{26-c}$, not primaries of
$\mathrm{Vir}_c$.  The two sequences differ: bar cohomology is
supported only at even weights and is $c$-independent; the
quasi-primary count is nonzero at odd weights and depends on the
specific gravity theory.

\textbf{(iii) The GRR holonomy level} recovers the monodromy of
$\cZ_{\mathrm{Vir}_c}$: the action of the mapping class group on
chiral homology.  For the Virasoro algebra, this is the
monodromy of the KZ connection, which encodes the
projective representation of $\mathrm{MCG}(\Sigma_g)$ on the
space of conformal blocks.  This is an invariant of the
\emph{algebra} (it is functorial in morphisms of chiral algebras,
not in changes of the module category or modular invariant).

\textbf{The obstruction is not analysis vs.\ algebra.}
It is \emph{algebra vs.\ category}: the bar complex sees
the chiral algebra $\mathrm{Vir}_c$; the gravitational
spectrum sees the choice of gravity theory built from
$\mathrm{Vir}_c$ (the specific multiplicities $\{m_h\}$
in the partition function, which are dynamical data external
to the algebra).

\subsection{Three routes toward the spectrum}

\paragraph{(a) Module bar construction.}
The module bar complex $\barB_{\mathrm{Vir}_c}(L(c,h))$ applied
to each irreducible module $L(c,h)$ produces a comodule over
$\barB(\mathrm{Vir}_c)$.  Its $K$-class carries the conformal
dimension~$h$.  Summing with the gravitational multiplicities
$\{m_h\}$:
\[
\sum_{h} m_h\,\mathrm{ch}\bigl(
R\pi_{1*}\,\barB_{\mathrm{Vir}_c}(L(c,h))\bigr)
\]
would recover the gravitational partition function at genus~$1$.
But this takes $\{m_h\}$ as \emph{input}: it computes the
partition function from the spectrum, rather than extracting
the spectrum from the bar complex.

\paragraph{(b) Conformal blocks as $D$-module.}
The conformal blocks local system on $\cM_{g,n}$ carries a flat
connection (the KZ/Hitchin connection) whose monodromy encodes
the modular $S$-matrix for rational theories.  At integrable
level for affine Kac--Moody algebras, the monodromy recovers the
primary dimensions via the Verlinde formula.

For the Virasoro algebra at generic~$c$, the conformal blocks
local system has irregular singularities and the representation
theory is non-semisimple (at irrational~$c$, all Verma modules
are irreducible and the spectrum is continuous).  Whether the
$D$-module structure of the conformal blocks sheaf on
$\overline{\cM}_{1,1}$ determines the gravitational spectrum
(after choosing a modular invariant) is a frontier question:
it amounts to recovering the spectral decomposition of a
function on $\mathrm{SL}(2,\Z)\backslash\mathfrak{H}$ from
the algebraic structure of its associated $D$-module.

\paragraph{(c) The Roelcke--Selberg step.}
The passage from the partition function
$Z(\tau,\bar\tau)$ to the constrained Epstein series
$\varepsilon^c_s$ uses the spectral decomposition of
$L^2(\mathrm{SL}(2,\Z)\backslash\mathfrak{H})$ into Eisenstein
series and Maass cusp forms.  This involves the hyperbolic
Laplacian, which is a real-analytic operator.
The scattering matrix involves $\zeta(2s)/\zeta(2s-1)$,
whose poles are the nontrivial zeros of the Riemann
zeta function.  No categorical or motivic construction is
known to produce this spectral decomposition.

The ``analytic'' step that is genuinely unavoidable is this
last one: from characters (which are algebraic for rational VOAs)
to spectral coefficients (which require harmonic analysis on
$\mathrm{SL}(2,\Z)\backslash\mathfrak{H}$).

\subsection{Honest status}

\begin{enumerate}[label=(\alph*)]
\item The bar complex defines a sheaf on $\overline{\cM}_g$
  with algebraic invariants ($R\Gamma$, $\mathrm{ch}$,
  $\mathrm{Hol}$, GRR).  The claim ``integration is analytic''
  was a false idea.

\item These algebraic invariants capture the algebra's modular
  structure: $\kappa$, bar cohomology dimensions, KZ/Hitchin
  monodromy.  They do \emph{not} capture the gravitational
  spectrum $\{m_h\}$ or the constrained Epstein
  $\varepsilon^c_s$.

\item The obstruction is \emph{algebra vs.\ category},
  not analysis vs.\ algebra.  Two gravity theories with the
  same Virasoro algebra but different spectra have the same
  bar complex.  This is a theorem, not a technical limitation.

\item The construction that does capture the spectrum---the
  conformal blocks $D$-module on $\cM_{g,n}$---is algebraic
  (Beilinson--Drinfeld) but requires the module category
  as input.  For rational theories, its monodromy ($S$-matrix)
  determines the primary dimensions.  For non-rational theories,
  the $D$-module has irregular singularities and the spectral
  question is at the frontier of geometric Langlands and
  irregular Hodge theory.

\item The single genuinely analytic step is the
  Roelcke--Selberg decomposition: from characters to spectral
  zeta.  This step is universal (it involves the Riemann
  zeta function, not the specific VOA) and has no known
  categorical substitute.

\item Status: the ``spectral question'' is \textbf{OPEN} in the
  sense that the algebraic $D$-module framework exists but
  has not been carried through for Virasoro at generic~$c$.
  The gap is not a missing construction but a missing
  computation: does the irregular $D$-module structure of the
  Virasoro conformal blocks local system on
  $\overline{\cM}_{1,1}$ determine the primary spectrum after
  choosing a modular invariant?
\end{enumerate}


% ============================================================
\section{Complex parameters in quantum gravity}
\label{sec:complex-parameters-qg}
% ============================================================

The gravitational bar complex of the Virasoro algebra produces
real spectral data: the modular characteristic $\kappa = c/2$,
the shadow tower $\{S_r(c)\}$, and the genus expansion
$F_g = \kappa\,\lambda_g^{\mathrm{FP}}$ are all rational
functions of the central charge~$c$.  The zeros of the
Riemann zeta function, which appear in the scattering matrix
of the Eisenstein spectral decomposition on $\cM_{1,1}$, live
at complex spectral parameters.  This section records the
definitive assessment of whether any extension of the bar
complex can access complex spectral data.

\subsection{The structural obstruction}

The obstruction is a theorem with three components.

\begin{enumerate}
\item The shadow tower on a single primary line is determined
by three genus-$0$ OPE invariants $(\kappa, \alpha, S_4)$
via Theorem~\ref{thm:riccati-algebraicity}: the generating
function $H(t) = t^2\sqrt{Q_L(t)}$ is algebraic of degree~$2$.
Three numbers constitute a finite-dimensional algebraic datum.

\item The zeta zeros $\{\rho_n = \tfrac{1}{2} + i\gamma_n\}$
constitute a countably infinite collection of complex numbers,
believed to be algebraically independent.  No finite-dimensional
algebraic structure determines such a set.

\item The Rankin--Selberg unfolding that connects the partition
function to the constrained Epstein $\varepsilon^c_s$ erases
the convolution-algebraic structure of the MC equation.  The
unfolded integral depends on Fourier coefficients alone, not on
the MC brackets.  The $L$-functions $L(s, f_j)$ in the
spectral decomposition are properties of the eigenforms~$f_j$,
independent of the algebra~$\cA$.
\end{enumerate}

\subsection{Eight candidate bridges: the census}

A complete investigation examined eight approaches to bridging
the gap from real bar-complex data to complex $L$-function
zeros.

\paragraph{Dead ends (five).}
Deligne categories, Wick rotation in~$c$, resonance/MC4
completion, Borel resummation, and spectral curves.
Each fails for a precise reason: Deligne categories
interpolate in a categorical variable, not a spectral one;
Wick rotation changes the shadow coefficients but not the
sewing lift $S_\cA(u)$, which is $c$-independent; MC4
completion solves a homological algebra problem, not a
number-theoretic one; Borel resummation requires factorial
divergence that the shadow series does not exhibit; and
spectral curves encode genus-$0$ OPE data, not the automorphic
spectrum.

\paragraph{Genuine mathematics, no bridge (two).}
The analytic continuation of the shadow tower and the analytic
Langlands programme produce real mathematical content but do
not bridge the real-to-complex gap.  The shadow tower
continuation produces quadratic number fields (the Galois
structure of the shadow extension $K_L/F$, with discriminant
$\mathrm{disc}(Q_L) = -320\,c^2/(5c{+}22)$ for Virasoro).
The analytic Langlands provides the correct framework for the
critical-level oper/shadow interpolation.  Both operate on the
function-field (geometric) side; the number-field (arithmetic)
descent remains the deepest unsolved problem.

\paragraph{The most compatible approach (one).}
The Connes--Marcolli non-commutative geometry programme has the
closest structural compatibility with the bar construction: the
Bost--Connes system has a Hecke algebra of symmetries
compatible with Verdier--Hecke commutation, and the Weil
positivity criterion is a trace-class condition on a specific
operator.  At the self-dual point $c = 13$, Verdier
self-duality provides an involution on the factorisation
coalgebra that could in principle feed into the non-commutative
trace formula.  In practice, the arithmetic descent (the bar
complex is defined over~$\mathbf{C}$, the Bost--Connes system
over~$\mathbf{Q}$) blocks the construction.

\subsection{What the gravitational bar complex provides}

The bar complex of the Virasoro algebra at central charge~$c$
provides five kinds of arithmetic data, all proved.

\begin{enumerate}
\item The exact genus expansion
$F_g = (c/2)\,\lambda_g^{\mathrm{FP}}$ at all genera.
\item The infinite shadow tower with growth rate
$\rho(c) = \sqrt{(36 + 2\Delta)/(c^2)}$ and discriminant
$\Delta = 80/(5c{+}22)$.
\item The four-class depth assignment (class~M for all
$c \neq 0$).
\item The quadratic number field
$K_L = \mathbf{Q}(\sqrt{\mathrm{disc}(Q_L)})$ at each
rational~$c$, with class numbers ranging from $h = 1$
(models $(9,10)$ and $(24,25)$) to $h = 3008$ (model $(28,29)$).
\item Route~C prime-locality: the MC recursion propagates
Hecke equivariance from low-arity character data to all
arities.
\end{enumerate}

None of these constrains the zeros of the Riemann zeta function
or any standard $L$-function.  The constraints are algebraic
(on OPE coefficients and their MC recursions); the zero
locations are analytic (determined by the Euler product and the
functional equation, not by the OPE).

\subsection{The gravitational dichotomy and the spectral divide}

The gravity dichotomy (critical string $c = 26$ versus
non-critical $c \neq 26$) is a statement about the
\emph{vanishing} of a specific class:
$\kappa_{\mathrm{eff}} = \kappa(\mathrm{matter})
+ \kappa(\mathrm{ghost}) = 0$ if and only if $c = 26$.
The self-dual point $c = 13$ is where the Koszul asymmetry
$\delta_\kappa = \kappa - \kappa' = 0$ vanishes.  These are
different conditions at different central charges (AP29).

The gravity dichotomy operates on the \emph{algebraic} axis:
whether the Clifford algebra $\mathrm{Cliff}(V, \lambda)$
with $\lambda = \kappa(\cB)\,\omega_g$ degenerates to an
exterior algebra ($\lambda = 0$ at $c = 26$) or remains a
matrix algebra ($\lambda \neq 0$ at $c \neq 26$).  This is a
statement about the transgression algebra, not about the
scattering matrix.  The spectral divide (real versus complex
spectral parameters) operates on the \emph{analytic} axis.
The two axes are orthogonal, and neither the gravitational
dichotomy nor the Koszul self-duality provides a bridge between
them.

\subsection{The Beilinson verdict for quantum gravity}

The following claims about the gravitational theory are
dismissed.

\begin{enumerate}
\item ``3d quantum gravity from Koszul duality constrains zeta
zeros.''  \emph{False.}  The Koszul duality of the Virasoro
algebra determines the complementary algebra $\mathrm{Vir}_{26-c}$
and the complementarity identity $\kappa + \kappa' = 13$.
These are algebraic constraints on OPE invariants, not analytic
constraints on $L$-function zeros.

\item ``The gravitational genus expansion encodes arithmetic
data beyond the character.''  \emph{Partially true, partially
false.}  The genus expansion $F_g = \kappa\,
\lambda_g^{\mathrm{FP}}$ is tautological (it lives in the
tautological ring of~$\overline{\cM}_g$).  The shadow tower
provides additional algebraic data (the quartic resonance
class, the discriminant, the depth classification).  But none
of this data accesses the primary spectrum or the partition
function, which are the objects that carry arithmetic content
through the Rankin--Selberg integral.

\item ``The gravitational partition function at $c = 26$ has
a distinguished relationship to the Riemann zeta function.''
\emph{False.}  The partition function of the
$c = 26$ Virasoro is $\eta(\tau)^{-24} = \Delta(\tau)^{-1}$,
the inverse of the Ramanujan discriminant form.  The
Rankin--Selberg $L$-function of $\Delta$ is
$L(s, \Delta \times \Delta) = \zeta(s{-}11)\,
L(s, \mathrm{Sym}^2\, \Delta)$, which has critical line
$\mathrm{Re}(s) = 23/2$, not $\mathrm{Re}(s) = 1/2$.
The zeta function $\zeta(s{-}11)$ appearing here is shifted
by~$11$ (the weight of~$\Delta$ minus~$1$).  The relationship
to the Riemann zeta is classical (Hecke theory), not a
consequence of the bar construction.
\end{enumerate}

\subsection{What survives for the gravitational theory}

The bar complex of the Virasoro algebra at central charge~$c$
is the algebraic engine of $3$d quantum gravity on
$\mathbf{C}_z \times \mathbf{R}_t$.  It determines:
the Swiss-cheese algebra structure (bar differential = chiral
product in~$\mathbf{C}$, bar coproduct = topological cutting
in~$\mathbf{R}$), the curved genus tower ($d^2 = \kappa\,
\omega_g$ at genus $g \ge 1$), the infinite shadow Postnikov
tower (all arities, all genera, proved by
Theorem~\ref{thm:mc2-bar-intrinsic}), the complementarity
decomposition (Theorem~C), and the completed platonic datum
$\Pi^{\mathrm{oc}}_X(\cA)$ packaging open/closed/modular data.

These are genuine results about the algebraic structure of the
gravitational theory.  They do not imply, and should not be
claimed to imply, any result about the analytic number theory
of $L$-functions.  The bar complex is the algebraic engine.
The primary spectrum is the analytic output.  The gap between
them is the gap between algebra and analysis, which no known
construction bridges.


% =================================================================
\section{The chiral Tamarkin problem}
\label{sec:chiral-tamarkin}
% =================================================================

A natural question about the bar construction is whether the
modular direction (the genus tower controlled by the
Getzler--Kapranov Feynman transform) is a genuine ``second
direction'' in the sense of Tamarkin's dimension-raising theorem.

The answer is: \emph{no and yes}.  The modular direction is not a
second holomorphic direction (it is operadic, not
factorisation-algebraic), but it is a genuine codimension-$1$
enhancement of the chiral
algebra, categorically independent of the Swiss-cheese centre,
and governed by its own universal construction.

\subsection{Two codimension-$1$ operations}

Tamarkin's theorem (2000): for an $\mathsf{E}_d$-algebra~$A$,
the Hochschild cochain complex $C^*(A,A)$ is an
$\mathsf{E}_{d+1}$-algebra.  The mechanism is convolution against
the little-discs cooperad.  The extra direction is topological.

For chiral algebras, there are two independent analogues:

\smallskip

\noindent
\textbf{(i) The Swiss-cheese centre} (the chiral Deligne--Tamarkin
theorem, proved in Volume~I).  The chiral Hochschild cochain complex
$C^\bullet_{\mathrm{ch}}(\cA,\cA)$ is the terminal local
$\SCchtop$-algebra.  This adds one topological direction~$\R$,
producing a $3$d HT theory on $X \times \R$.

\smallskip

\noindent
\textbf{(ii) The modular convolution} (Getzler--Kapranov, proved in
Volume~I).  The modular convolution dg Lie algebra $\gAmod$ extends
genus-$0$ data to all genera.  The mechanism is convolution against
the modular cooperad of stable curves.  The BV operator
$\hbar\Delta$ (non-separating clutching
$\overline\cM_{g,n+2} \to \overline\cM_{g+1,n}$) raises genus
by~$1$.

\smallskip

\noindent
These are different functors: (i) uses the Swiss-cheese coloured
operad (configuration spaces of discs in $\C \times \R$);
(ii) uses the modular operad (stable curves of all genera with
clutching maps).  Both are convolutions against cooperads, and
both produce Maurer--Cartan-controlling Lie-type algebras.

The open/closed MC element $\Theta^{\mathrm{oc}}_\cA = \Theta_\cA
+ \sum_j \mu^{\cM_j}$ packages both into a single algebraic
object.  Its closed projection is $\Theta_\cA$ (the modular
convolution MC element).  Its open projection gives the $\Ainf$
boundary module structure.  Its mixed projection gives the
bulk-boundary coupling (the chiral derived centre acting on the
boundary via the Swiss-cheese structure).  Its clutching sector
gives the genus tower.


\subsection{The modular open/closed cooperad conjecture}

The hierarchy of Tamarkin-type theorems is:

\begin{center}
\renewcommand{\arraystretch}{1.4}
\begin{tabular}{@{}p{3.3cm}p{3.3cm}p{5.5cm}@{}}
\textsc{Tamarkin:}
  & $A$ is $\mathsf{E}_d$
  & $\Longrightarrow\; C^*(A,A)$ is $\mathsf{E}_{d+1}$ \\[2pt]
\textsc{Chiral DT:}
  & $\cA$ is chiral
  & $\Longrightarrow\; (C^\bullet_{\mathrm{ch}}(\cA,\cA), \cA)$
    is $\SCchtop$-algebra \\[2pt]
\textsc{CT-$2$:}
  & $\cA$ chirally Koszul
  & $\Longrightarrow\;$ Swiss-cheese extends to
    modular Swiss-cheese
\end{tabular}
\end{center}

\begin{conjecture}[CT-$2$: the modular open/closed cooperad]
\label{conj:ct2-vol2}
The open-sector factorisation dg-category $\cC_{\mathrm{op}}$
admits a modular cooperad structure whose cocomposition maps are
compatible with clutching on bordered Fulton--MacPherson
compactifications, whose trace recovers the annulus identification
$\HH_*(\cA_b) \simeq \Tr_\cA$ at genus~$1$, and whose full
structure determines and is determined by the open/closed MC
element~$\Theta^{\mathrm{oc}}_\cA$.
\end{conjecture}

\noindent\textbf{Status.}\quad
CT-$2$ is a research programme with at least six independent
sub-problems: (P1)~dualizability of~$\cC_{\mathrm{op}}$
in~$\mathrm{dgCat}$; (P2)~annular bordered FM compactification
(the cyclic case is not yet constructed---the genus-$1$ annulus
trace is proved by Ayala--Francis excision, not by bordered FM);
(P3)~construction of functorial cocomposition maps along bordered
FM clutching for all stable topological types;
(P4)~codimension-$2$ coassociativity (the MC equation provides
codimension-$1$ identities; coassociativity is codimension~$2$);
(P5)~$(\infty,2)$-categorical lift from chain-level cooperad to
cooperad in dg-categories; (P6)~coderived compatibility at genus
$g \geq 1$ (curved bar forces passage to coderived categories).
The MC equation for $\Theta^{\mathrm{oc}}$ is proved at all
genera.  The genus-$0$ part is the chiral Deligne--Tamarkin
theorem (proved).  This is Stage~$4$ of the four-stage
architecture.

\begin{remark}[Why this is the correct formulation]
The conjecture is not ``the shadow tower is a chiral algebra on
moduli'' (that is ill-defined: $\overline\cM_g$ has dimension
$3g{-}3$, not~$1$).  The conjecture is that the open-sector
category carries a modular cooperad structure compatible with
the proved MC equation.  The modular cooperad is the algebraic
structure that governs the genus direction; it is not a
factorisation algebra on a Ran space, but an operadic structure
indexed by boundary strata of moduli spaces.  The distinction
between factorisation (points colliding on a fixed curve) and
composition (curves degenerating along boundary divisors) is
categorical.
\end{remark}


\subsection{The BCOV parallel in the HT setting}
\label{subsec:bcov-ht}

The genus-by-genus MC hierarchy from Volume~I,
\[
d_{\Theta^{(0)}}(\Theta^{(g)})
= -\tfrac{1}{2}\!\sum_{g_1+g_2=g}
  [\Theta^{(g_1)}, \Theta^{(g_2)}]
- \Delta(\Theta^{(g-1)})
- d_{\mathrm{sew}}^{(g)} - d_{\mathrm{pf}}^{(g)},
\]
has the structure of the BCOV holomorphic anomaly equation:
both are genus recursions governed by the modular operad, with a
bilinear term from separating degenerations and a genus-raising
term from non-separating degenerations.

In the $3$d HT setting of this volume, the parallel is sharper.
The bulk algebra $\Zder(\cA) = C^\bullet_{\mathrm{ch}}(\cA_b,
\cA_b)$ is the universal closed-sector algebra
(Theorem~\ref{thm:thqg-swiss-cheese}).  At genus~$0$, the bulk is
the chiral derived centre.  At genus $g \geq 1$, the curved
Swiss-cheese structure $d^2 = \kappa \cdot \omega_g$ infects the
bulk via the Gauss--Manin connection.  The total differential
$D_g = d_{\mathrm{fib}} + \nabla^{\mathrm{GM}}_g$ restores
flatness by absorbing the curvature into the connection.

\begin{warning}[The BCOV analogy is combinatorial, not structural]
\label{warn:bcov-not-structural-vol2}
The MC hierarchy and the BCOV holomorphic anomaly equation share a
combinatorial backbone (genus recursion from the modular operad),
but they are not structurally isomorphic.  The shadow tower has no
holomorphic anomaly ($\bar\partial = 0$ on algebraic classes), no
non-canonical propagator (the bar propagator $d\log E$ is
canonical), and two extra terms ($d_{\mathrm{sew}}$,
$d_{\mathrm{pf}}$) with no BCOV counterpart.  The correct physics
analogue is the DVV/Virasoro recursion for CohFTs, via Teleman
reconstruction and the Eynard--Orantin identification.
\end{warning}

\noindent
What does survive: the Gauss--Manin variation on $1$D deformation
families.  For $\{\cA_s\}_{s \in B}$, the shadow generating
function $H(t;s) = t^2\sqrt{Q_L(t;s)}$ satisfies a first-order
ODE in~$s$ controlled by $\kappa'(s)$, $\alpha'(s)$, $S_4'(s)$.
The flat modular connection $D_\Theta = d + [\Theta_\cA, -]$ on
bar cohomology over $\cM_g$ gives the all-genera extension.

The decisive asymmetry: the shadow tower has \emph{no holomorphic
anomaly} ($\Theta^{(g)}$ is algebraic, so $\bar\partial = 0$
automatically), and the shadow partition function \emph{converges}
(Bernoulli decay $1/(2\pi)^{2g}$, versus $(2g)!$ divergence for
BCOV).  The algebraic bar construction is the universal structure;
BCOV is one analytic incarnation.


\subsection{Chromatic structure of the bar complex}
\label{subsec:chromatic-bar}

The scalar shadow tower is at chromatic height~$0$ (additive
formal group, $\hat{A}$-genus).  But the bar complex on curves
of different type uses propagators of different height:
$1/(z{-}w)$ on $\bP^1$ (rational, height~$\infty$),
$\pi\cot(\pi z)$ on $\C^*$ (trigonometric, height~$1$),
$\vartheta_1'/\vartheta_1$ on $E_\tau$ (elliptic, height~$2$).

The rational/trigonometric/elliptic hierarchy of quantum groups
IS the chromatic tower.  The Fay identity on $E_\tau$ is the
genus-$1$ MC equation with elliptic propagator.

\begin{conjecture}[Chromatic bar--cobar duality]
\label{conj:chromatic-vol2}
The bar--cobar adjunction on a family of elliptic curves
$E \to S$ defines a sheaf of $\Etwo$-algebras on~$S$ whose
associated formal group is the universal elliptic formal group.
The resulting Koszul duality theory determines a
$\mathrm{TMF}$-module structure on factorisation coalgebras
on $E/S$.
\end{conjecture}

\noindent
The genus spectral sequence---with $E_1$ page isolating tree
($p{=}0$), one-loop ($p{=}1$), genus-$2$ shell ($p{=}2$)
data---admits a conjectural identification with the chromatic
filtration.  Whether the shadow depth classification G/L/C/M
corresponds to chromatic complexity is a speculative question at
the interface of modular Koszul duality and stable homotopy theory.

\sep

The three conjectures of this section
(CT-$2$, the BCOV reformulation, and the chromatic lift)
are not independent.  The modular cooperad conjecture CT-$2$
would provide the categorical infrastructure in which both the
BCOV equations and the chromatic structure would live.
The BCOV reformulation requires understanding how the
shadow connection $\nabla^{\mathrm{sh}}$ extends from
$1$-dimensional deformation slices to the full moduli of chiral
algebras---and this extension is controlled by the cooperad
cocomposition maps of CT-$2$.
The chromatic conjecture requires showing that the elliptic bar
complex varies correctly over the modular curve---and the
relevant variation is a crystal of factorisation coalgebras,
whose coherence is governed by the same modular cooperad.

The unifying perspective: the MC element $\Theta_\cA$ in
$\gAmod$ simultaneously encodes the spatial dimension-raising
(its genus-$0$ binary shadow $r(z)$ is the R-matrix, providing
the $\mathsf{E}_1$/Yangian structure in the topological
direction) and the genus dimension-raising (its higher-genus
components are the quantum corrections, living on
$\overline\cM_g$).  Quantisation and genus-raising are two
faces of the same MC element, expressed in different projections
of the shadow Postnikov tower.


% =================================================================
\section{The bar construction as string field theory}
\label{sec:bar-sft}
% =================================================================

The bar construction is a string field theory functor.  The
$A_\infty$ operations from the chiral Hochschild complex are
Zwiebach's closed string vertices ($m_1 = Q_{\mathrm{BRST}}$,
$m_2 =$ string product, $m_3 =$ four-string vertex, higher
$m_n =$ contact terms), and the action is
$S[\Psi] = \frac{1}{2}\langle\Psi, Q\Psi\rangle
+ \sum_{n \geq 3} \frac{1}{n!}
  \langle\Psi, m_n(\Psi,\ldots,\Psi)\rangle$.
This is Zwiebach's theorem (1993), restated as
Theorem~\ref{thm:string-field-theory-hochschild} in Volume~I.

The genus-$0$ BV/bar identification is proved.  The full
chain-level identification at genus $g \geq 1$ is heuristic
(the MC equation is proved; the physical identification
BV $=$ bar is Conjecture~\ref{conj:master-bv-brst}).

\subsection{Four objects and their string-theoretic roles}

The open/closed string field theory of a chiral algebra $\cA$
on $\C \times \bR$ involves four objects:
\begin{itemize}[nosep]
\item $\cA$ (boundary algebra) $=$ open-string sector.
\item $\barB(\cA)$ (bar coalgebra) $=$ classifies twisting
  morphisms (universal couplings between $\cA$ and $\cA^!$).
  This is the open-string configuration space, not the
  closed-string state space.
\item $\Zder(\cA) = C^\bullet_{\mathrm{ch}}(\cA,\cA)$
  (derived centre) $=$ universal bulk algebra (closed-string
  observables).
\item $\Theta^{\mathrm{oc}}_\cA$ (open/closed MC element) $=$
  the full open/closed string field, packaging all four sectors.
\end{itemize}
The bar complex and the derived centre are \emph{different}
objects playing \emph{different} roles:
bar $=$ couplings; centre $=$ observables.
Open/closed string duality is the interplay between them, not
the identification of one with the other.

\subsection{Kodaira--Spencer gravity and open/closed duality}

Volume~II's worked example (Section~\ref{subsec:KS-defect})
identifies the large-$N$ BRST cohomology of the gauged
$\beta\gamma$ system with the Kodaira--Spencer chiral algebra
of BCOV.  The passage from open-string to closed-string
degrees of freedom is controlled by \emph{Koszul duality}
($\cA_{\mathrm{open}} \mapsto \cA^!_{\mathrm{open}}$ via
Verdier duality on bar cohomology) and the \emph{derived
centre} $C^\bullet_{\mathrm{ch}}(\cA_{\mathrm{open}},
\cA_{\mathrm{open}})$ as the universal bulk.  Bar-cobar
inversion $\Omega B(\cA_{\mathrm{open}}) \simeq
\cA_{\mathrm{open}}$ recovers the original boundary algebra
(Theorem~B), not the KS bulk~(AP34).

The structural parallel between BCOV and the modular MC
hierarchy arises because both are string field theories governed
by the modular operad.  The BCOV system is the analytic
incarnation (metric-dependent propagator, holomorphic anomaly,
divergent).  The shadow tower is the algebraic incarnation
(canonical propagator, no anomaly, convergent).

\begin{question}[Open/closed duality: bar vs derived centre]
\label{q:open-closed-vol2}
The bar complex classifies twisting morphisms.  The derived
centre classifies bulk observables.  The Swiss-cheese structure
gives the local answer: $(\Zder(\cA), \cA)$ is the terminal
$\SCchtop$-pair.  The modular extension gives
$\Theta^{\mathrm{oc}}_\cA$.  But what is the precise categorical
relationship between the category of twisting morphisms
(controlled by $\barB(\cA)$) and the category of bulk operators
(controlled by $\Zder(\cA)$)?

In the topological setting, the Hochschild cochains of an
$\Eone$-algebra are computed by the bar complex: $C^*(A,A)
\simeq \RHom_{A \otimes A^{\mathrm{op}}}(A,A) \simeq
\RHom(\barB(A), A)$.  In the chiral setting, the chiral
Hochschild cochains are $C^\bullet_{\mathrm{ch}}(\cA,\cA)$,
which is the operadic centre, \emph{not} $\RHom(\barB(\cA),
\cA)$.  The gap between the operadic centre (the derived
centre) and the convolution $\Hom$ from the bar coalgebra
(the twisting morphism space) is where open/closed duality
lives.  Closing this gap is the mathematical content of
open/closed string duality in the chiral setting.
\end{question}


\end{document}
